\chapter{Uniqueness for degenerate equations}
\label{chap:degenerate-uniqueness}

Throughout this chapter, $X$ is a connected compact complex manifold of dimension $n\geq1$, $\omega$ is a Hermitian form and $\mathrm{d}V$ is a smooth positive volume form on $X$. As before, $\omega_u=\omega+\ddc u$, and for bounded $\omega$-plurisubharmonic $u$ we write $\mathrm{MA}_\omega(u)=\omega_u^n$ for the Bedford--Taylor product of \cref{prop:bounded-products}. The preceding chapter shows that, for a nonnegative nonzero density $f\in L^p(X,\mathrm{d}V)$ with $p>1$, the constant $c$ in the equation $\omega_u^n=cf\,\mathrm{d}V$ is determined by $f$. We now ask when the potential is unique up to an additive constant.

For real potentials $u,v$, with the regularity stated in each result, set
\[
 A=\omega+\ddc u,\qquad B=\omega+\ddc v,\qquad w=u-v,
 \qquad T=\sum_{j=0}^{n-1}A^j\wedge B^{n-1-j}.
\]
Then
\begin{equation}\label{uni-eq:linearization}
 \ddc w\wedge T=A^n-B^n.
\end{equation}
Thus, if $A^n=B^n$, the difference $w$ solves the linear equation $\ddc w\wedge T=0$. When $A,B\geq0$, the polarized cofactor $T$ is semipositive, but it may be degenerate and need not be closed. For Sobolev or $C^{1,1}$ potentials, products of coefficients in these identities are understood almost everywhere.

For $C^{1,1}$ solutions with the same continuous density, strict positivity of the density gives local uniform ellipticity of this linear equation. The main issue is to locate the extrema of $w$ in these elliptic regions and to decide whether their values propagate across the zero set of the density. A positive weight making $T$ $\ddc$-closed gives another route: a global energy identity forces the derivative of $w$ to vanish in the directions of the cofactor. For bounded potentials, an energy measure replaces this pointwise calculation.

If the background has invariant mass, ordinary comparison proves uniqueness of bounded solutions for $L^p$ densities, $p>1$, without any connectedness condition on the set where the density is positive. For a general Hermitian background, however, two hypersurfaces on which the density vanishes can separate regions in which $w$ attains different extremal values. The examples on $\mathbb P^3$ at the end of the chapter first place this freedom in an open band where the density vanishes, and then keep it with a density having only quadratic zeros, which can even be chosen real analytic. They show that smallness of the zero set alone does not restore uniqueness.

\section{Extrema and measurable ellipticity}
\label{uni-sec:anchors}

Let $X$ be a connected compact complex manifold of dimension $n\geq1$, let $\omega$ be a Hermitian form on $X$, and fix a smooth positive volume form $\mathrm{d}V$. For real potentials $u,v$ of the regularity specified below, put $A=\omega+\ddc u$, $B=\omega+\ddc v$, $w=u-v$ and $T=\sum_{j=0}^{n-1}A^j\wedge B^{n-1-j}$. Products of measurable coefficients are interpreted almost everywhere. Suppose $u,v\in C^{1,1}(X,\mathbb R)$, $A,B\geq0$ almost everywhere, and
\begin{equation}
  A^n=B^n=f\,\mathrm{d}V,\qquad f\in C^0(X),\qquad f\geq0,
  \label{uni-eq:density}
\end{equation}
where the coefficient products are understood almost everywhere.  Fix a
smooth positive $(1,1)$-form $\beta$, put $H=A+B$, and set
\begin{equation}
 E_{\rm unif}
 =\left\{x:\begin{array}{l}
 \text{there are a neighborhood }V\ni x\text{ and }c>0\\
 \text{with }H\geq c\beta\text{ almost everywhere on }V
 \end{array}\right\}.
 \label{uni-eq:Eunif}
\end{equation}
This open set is independent of $\beta$ and of the representatives of the
measurable coefficients.

\begin{theorem}\label{uni-thm:anchors}
Let $u,v\in C^{1,1}(X)$, assume that $A,B\geq0$ almost everywhere, and
assume \eqref{uni-eq:density}.  If $\chi\in C^2(X,\mathbb R)$ satisfies
$C:=\omega+\ddc\chi>0$, put
\[
  P_\chi=\{x\in X:f(x)\mathrm{d}V_x\geq C_x^n\},\qquad U=\{f>0\}.
\]
Then $P_\chi$ is a nonempty compact subset of $U$, and
$U\subset E_{\rm unif}$.  There are components $G_+,G_-$ of
$E_{\rm unif}$, each meeting $P_\chi$, such that
\begin{equation}
 w\equiv\max_Xw\quad\hbox{on }\overline{G_+},\qquad
 w\equiv\min_Xw\quad\hbox{on }\overline{G_-}.
 \label{uni-eq:extrema-components}
\end{equation}
If $w$ is nonconstant, these closures are disjoint.  Consequently $w$ is
constant if any one of the following holds:
\begin{enumerate}
\item every two components of $E_{\rm unif}$ meeting $P_\chi$ have
      intersecting closures;
\item $E_{\rm unif}$ is connected;
\item $E_{\rm unif}$ is dense and has at most two components.
\end{enumerate}
The same three conclusions hold with $U$ in place of $E_{\rm unif}$.
In particular, uniqueness holds when $\{f=0\}$ is a nonempty connected
embedded $C^1$ real hypersurface without boundary.  No orientability or
coorientability assumption is needed.
\end{theorem}

We begin with two local lemmas.

\begin{lemma}\label{uni-lem:goodcontact}
In Euclidean space $\mathbb R^m$, let $B(p,r)$ be the ball of center $p$ and radius $r>0$, and let $I$ be the identity matrix. Let $g\in C^{1,1}$ near $\overline{B(p,r)}\subset\mathbb R^m$, suppose
$g\leq g(p)$ on the ball, and let $E$ be any full-measure subset on
which $g$ has a second-order Taylor expansion.  For every $\delta>0$ and
$0<\rho<r$, there is $x\in E\cap B(p,\rho)$ such that
\[
    D^2g(x)\leq2\delta I.
\]
\end{lemma}

\begin{proof}
Put $g_\delta(x)=g(x)-\delta|x-p|^2$.  The gradient map
$F=\nabla g_\delta$ is Lipschitz on the ball, so it maps the null set
$N=\overline{B(p,r)}\setminus E$ to a null set.  Choose
$s\in B(0,\delta\rho)\setminus F(N)$ and maximize
$g_\delta(x)-s\cdot(x-p)$ on the closed ball.  At any point with
$|x-p|\geq\rho$, its value is strictly smaller than its value at $p$,
because
\[
 g(x)-g(p)-\delta|x-p|^2-s\cdot(x-p)
 \leq-\delta|x-p|^2+|s|\,|x-p|<0.
\]
Thus a maximizer $x$ lies in $B(p,\rho)$.  Its first derivative gives
$F(x)=s$, so $x\notin N$.  The second-order Taylor expansion at this
upper contact point gives $D^2g(x)\leq2\delta I$.
The null-set assertion for a locally Lipschitz gradient follows by a
countable cover by sets on which it is Lipschitz.  The full-measure set of
second-order Taylor points exists since the first derivatives of a
$C^{1,1}$ function are Lipschitz.
\end{proof}

\begin{lemma}\label{uni-lem:measSMP}
Let $a=(a^{ij})$ be a measurable symmetric matrix field on a connected
domain $\Omega\subset\mathbb R^m$. For each compact $K\Subset\Omega$, suppose there are constants $0<\lambda_K\leq\Lambda_K<\infty$ such that
\[
   \lambda_K I\leq a(x)\leq\Lambda_KI
\]
almost everywhere on $K$, where $I$ is the identity matrix. Repeated indices $i,j=1,\ldots,m$ are summed, and $z_{ij}=\partial_i\partial_jz$. If
$z\in W^{2,\infty}_{\rm loc}(\Omega)$,
$a^{ij}z_{ij}\geq0$ almost everywhere, and $z$ attains its supremum at an
interior point, then $z$ is constant.
\end{lemma}

\begin{proof}
We first prove strict weak comparison on a relatively compact domain:
if $b\in C^{1,1}$, $Lb:=a^{ij}b_{ij}\geq\kappa>0$ almost everywhere,
and $b\leq M$ on the boundary, then $b\leq M$ inside.  Otherwise an
interior maximum lies in a closed coordinate ball on which
$a\leq\Lambda I$.  Apply Lemma~\ref{uni-lem:goodcontact} to the full-measure
set where both the differential inequality and the Taylor expansion hold.
Choosing $2\delta m\Lambda<\kappa$ gives the contradiction
$\kappa\leq Lb\leq2\delta\operatorname{tr}a\leq2\delta m\Lambda$.

Take the $C^{1,1}$ representative of $z$, write $M=\sup_\Omega z$,
and suppose $K=\{z=M\}$ is nonempty and proper.  Near a relative boundary
point of $K$, choose $y\notin K$ so that the ball of radius
$r=\operatorname{dist}(y,K)>0$ is compactly contained in $\Omega$ and
touches $K$ at $x_0$.  On the closed annulus
$r/2\leq|x-y|\leq r$, use fixed bounds $\lambda I\leq a\leq\Lambda I$
and set
\[
 b(x)=\mathrm{e}^{-\alpha|x-y|^2}-\mathrm{e}^{-\alpha r^2}.
\]
For sufficiently large $\alpha$,
\[
 Lb=\mathrm{e}^{-\alpha|x-y|^2}
       \bigl(4\alpha^2 a^{ij}(x_i-y_i)(x_j-y_j)
                       -2\alpha\operatorname{tr}a\bigr)\geq\kappa>0
\]
almost everywhere on the annulus.  On its inner sphere $z<M$ uniformly;
on its outer sphere $b=0$.  Choose $\eta>0$ small enough that
$z+\eta b\leq M$ on both spheres.  Strict weak comparison gives this
inequality throughout the annulus.  At $x_0$, however, $\nabla z=0$
and $b$ increases to first order in the inward radial direction, which
contradicts that inequality.  Thus $K$ has no relative boundary and
connectedness proves the assertion.  This argument uses no derivative of
the measurable coefficients.
\end{proof}

At a point where $A,B$ are semipositive, define the cofactor quadratic form on real cotangent vectors
\begin{equation}
 Q_x(\xi)\mathrm{d}V_x
   =\mathrm{i}\,\xi^{1,0}\wedge\overline{\xi^{1,0}}\wedge T_x .
 \label{uni-eq:Q}
\end{equation}
Mixed products of semipositive $(1,1)$-forms are semipositive in this
covector-testing order.  Expanding $H^{n-1}$, $H=A+B$, gives
\begin{equation}
 2^{-(n-1)}H^{n-1}\leq T\leq H^{n-1}.
 \label{uni-eq:cofactor-comparison}
\end{equation}
At a point, choose a $\beta$-unitary coframe that diagonalizes $H$, and write $\lambda_1,\ldots,\lambda_n$ for its eigenvalues. In that coframe the
coefficient of
$\mathrm{i}\alpha\wedge\bar\alpha\wedge H^{n-1}$, for
$\alpha=\sum_{k=1}^n b_k\,\mathrm{d}z_k$, with complex coefficients $b_k$, is a positive dimensional constant times
\begin{equation}
  \sum_{k=1}^n |b_k|^2\prod_{\ell\ne k}\lambda_\ell.
 \label{uni-eq:cofactor-eigen}
\end{equation}
For $n\geq2$, \eqref{uni-eq:cofactor-comparison}--\eqref{uni-eq:cofactor-eigen}
show that $Q$ is locally uniformly elliptic exactly where $H$ has a
local essential positive lower bound.  The reverse implication uses the
local upper bound for $H$: testing each eigendirection bounds every
$\prod_{\ell\ne k}\lambda_\ell$ below, and choosing $k$ different from
an index of a smallest eigenvalue bounds that eigenvalue below.  For $n=1$
one has $T=1$, so the linearized equation is elliptic everywhere.

\begin{proof}[of Theorem~\ref{uni-thm:anchors}]
We first locate the extrema.  Fix $s>0$ and maximize
\begin{equation}
  w-s(v-\chi).
 \label{uni-eq:pertmax}
\end{equation}
For smooth data, the complex Hessian inequality at a maximizer is
\begin{equation}
  (1+s)B\geq A+sC.
 \label{uni-eq:matrixcontact}
\end{equation}
The positive term $sC$ first gives $B>0$, and determinant equality then
gives $A>0$.  In a $B$-unitary frame put
$J=B^{-1/2}AB^{-1/2}$, $K=B^{-1/2}CB^{-1/2}$.  Since $\det J=1$,
arithmetic--geometric mean and \eqref{uni-eq:matrixcontact} give
\[
 n+s\operatorname{tr}K\leq\operatorname{tr}J+s\operatorname{tr}K
 \leq n(1+s).
\]
Consequently $\operatorname{tr}K\leq n$, $B\geq C/n$, and
$\det K\leq(\operatorname{tr}K/n)^n\leq1$.  In particular $B^n\geq C^n$.

For $C^{1,1}$ potentials, fix a maximizer $p$ of
\eqref{uni-eq:pertmax} and apply Lemma~\ref{uni-lem:goodcontact} on a coordinate
ball, choosing radii tending to zero and errors $\delta_j\downarrow0$.
Include in its full-measure set the points where both Hessians, their
semipositivity, and their common determinant are defined.  With
$\beta_0=\ddc|z|^2$, the resulting points $x_j\to p$ satisfy
\[
 A(x_j)+s C_j(x_j)\leq(1+s)B(x_j),\qquad
 C_j=C-(\delta_j/s)\beta_0>0
\]
for large $j$.  The matrix argument just given yields
$f(x_j)\mathrm{d}V_{x_j}\geq C_j(x_j)^n$.  Continuity of $f,\mathrm{d}V,C$, not
continuity of the Hessians, gives $p\in P_\chi$.  Thus every maximizer of
the perturbed function lies in $P_\chi$.  Letting $s\downarrow0$
places an actual maximizer of $w$ in $P_\chi$.  Applying the same argument
to $-w$ places a minimizer there as well.  Compactness and positivity of
$C$ show that $P_\chi$ is compact, nonempty, and contained in $U$.

On a relatively compact ball in $U$, the determinants of $A$ and $B$
have a positive lower bound, while their coefficients have an essential
upper bound.  Every eigenvalue is therefore bounded below by the determinant
lower bound divided by a fixed power of the upper bound.  Thus
$U\subset E_{\rm unif}$.

By \eqref{uni-eq:linearization}, $w$ solves a measurable nondivergence equation
whose principal quadratic form is a fixed positive multiple of $Q$.  On
each relatively compact coordinate ball in $E_{\rm unif}$, its coefficients
are measurable and uniformly elliptic.  Lemma~\ref{uni-lem:measSMP} makes $w$
constant on the component containing its maximum found above, and makes $-w$
constant on the component containing its minimum found above.  Continuity
extends the two constant values to the closures, proving
\eqref{uni-eq:extrema-components}.  Distinct extrema force disjoint closures.

The same argument applies on $U$ itself: it contains the extremal points in $P_\chi$, and the cofactor is locally uniformly elliptic there. This proves the assertions with $U$ in place of $E_{\rm unif}$.

The component criteria are immediate except for the dense two-component
case.  If an open dense subset of connected $X$ has at most two components
with disjoint closures, those closures are disjoint nonempty closed sets
whose union is $X$, contradicting connectedness.  Finally, when
$Z=\{f=0\}$ is a connected embedded real hypersurface, flattening charts
show that all local sides belong to one fixed unordered pair of components
of $X\setminus Z$.  Hence $U$ has at most two components; it is dense
because a hypersurface has empty interior.  This proves the last assertion.
\end{proof}


\section{A square-root threshold for density sublevels}
\label{uni-sec:threshold}

The topology of the positive-density set can be controlled by the volume of small sublevels. A nonnegative $C^{1,1}$ density grows at most quadratically away from a zero. If its zero set separates a small ball, a family of transverse segments therefore contributes at least a constant times $\sqrt t$ to the volume of $\{f\leq t\}$. Excluding that lower bound gives connectedness and hence uniqueness.

Write $\operatorname{Vol}(E)=\int_E\mathrm{d}V$ for Borel sets $E$.

\begin{theorem}\label{uni-thm:threshold}
Let $u,v\in C^2(X)$, $A,B\geq0$, and
$A^n=B^n=f\,\mathrm{d}V$, where $f\geq0$ is $C^{1,1}$.  If
\begin{equation}
  \liminf_{t\downarrow0}
  \frac{\operatorname{Vol}\{f\leq t\}}{\sqrt t}=0,
  \label{uni-eq:sqrtcriterion}
\end{equation}
then $\{f>0\}$ is connected and $u-v$ is constant.  In particular this
holds when $f^{-1/2}\in L^1(\mathrm{d}V)$, with $f^{-1/2}=+\infty$ on $\{f=0\}$.
The exponent $1/2$ is sharp in every dimension $n\geq3$.
\end{theorem}

\begin{proof}
Write $Z=\{f=0\}$.  Condition \eqref{uni-eq:sqrtcriterion} first implies
$\operatorname{Vol}(Z)=0$.  We show that the complement of $Z$ in every sufficiently
small coordinate ball is connected.  Fix such a ball with convex Euclidean
image.  If its complement were disconnected, choose points $a,b$ in
different components and small closed endpoint disks perpendicular to the
segment $[a,b]$, contained in those respective components.  By making
the disks small, all parallel segments between corresponding points stay
in the ball.  Each segment meets $Z$, since otherwise it would connect
the two components.

On a slightly smaller fixed coordinate region the gradient of $f$ is
Lipschitz with some constant $K>0$.  Nonnegativity gives
$\nabla f(z)=0$ at every zero, hence
\begin{equation}\label{uni-eq:quadratic-upper}
 f(x)\leq\frac K2|x-z|^2
\end{equation}
whenever the segment from $z$ to $x$ stays in that region.  The endpoint
disks are compact and disjoint from $Z$, so their nearby segment portions
are uniformly free of zeros.  Thus every parallel segment contains a zero
a fixed positive distance from its endpoints.  For small $t>0$,
\eqref{uni-eq:quadratic-upper} supplies on that segment an interval of length
$2\sqrt{2t/K}$ belonging to $\{f\leq t\}$.  Integration over the
positive-measure transverse disk, and comparability of $\mathrm{d}V$ with
coordinate volume, give
\[
 \operatorname{Vol}\{f\leq t\}\geq c\sqrt t
\]
for all sufficiently small $t$, contradicting
\eqref{uni-eq:sqrtcriterion}.  This proves the local connectedness assertion.

Each such ball has a nonempty complement of $Z$, since $Z$ has zero
volume, and that complement belongs to one component of $U=X\setminus Z$.
Two overlapping balls have the same component label: their open overlap
contains a point outside $Z$.  These locally constant labels on
connected $X$ are therefore all the same.  Hence $U$ is connected,
and Theorem~\ref{uni-thm:anchors} gives uniqueness.

If $f^{-1/2}\in L^1$, then, outside the null zero set,
\[
 \frac{\operatorname{Vol}\{f\leq t\}}{\sqrt t}
 \leq\int_{\{f\leq t\}}f^{-1/2}\,\mathrm{d}V\longrightarrow0
\]
by dominated convergence.  This proves the stated sufficient integrability
condition.  The sharpness assertion follows from the examples in
Section~\ref{uni-sec:examples}, as quantified at the end of that section.
\end{proof}

\begin{remark}
The regularity in this criterion is a condition on the \emph{density}:
the quadratic upper bound at its zeros is the decisive input.
It is different from the $C^{1,1}$ condition on the \emph{potentials}
in the preceding uniqueness theorem.
\end{remark}


\section{Polarized cofactors and propagation}
\label{uni-sec:cofactor}

Uniform ellipticity is sufficient for propagation of an extremum, but it is not necessary. The range of the cofactor may already connect enough points to force constancy. For $C^2$ potentials, the quadratic form $Q$ in \eqref{uni-eq:Q} is continuous. Its polarization defines $Q^\sharp:T^*X\to TX$ by $\eta(Q^\sharp\xi)=Q(\xi,\eta)$. A reachable set from a point consists of endpoints of absolutely continuous curves starting there with velocity in $\operatorname{im}Q^\sharp$ almost everywhere. A control orbit consists of endpoints of finite concatenations of flows of locally Lipschitz vector fields taking values in that range.

A positive weight satisfying $\ddc(hT)=0$ gives an energy identity controlling $\mathrm{d}w$ everywhere. Without such a weight, a degenerate maximum principle gives propagation only from the extremal sets. This distinction accounts for the different density assumptions on reachable sets and orbits in the theorem.

\begin{theorem}\label{uni-thm:cofactor-intro}
Let $u,v\in C^2(X)$, $A,B\geq0$, and $A^n=B^n$.
\begin{enumerate}
\item If there is a continuous $h>0$ such that
$\ddc(hT)=0$ in distributions, then
\begin{equation}
  \mathrm{i}\,\partial w\wedge\bar\partial w\wedge T\equiv0.
  \label{uni-eq:energyzero}
\end{equation}
It follows that $w$ is constant if the strict positivity locus of the
cofactor is dense, or if one cofactor control-reachable set is dense.
\item Without the weight hypothesis, the global maximum and minimum sets of
$w$ are invariant under every locally Lipschitz vector field taking values
in the range of the cofactor symbol.  If every such control orbit is dense,
then $w$ is constant.  For smooth data it is enough that smooth range
fields, for example the cofactor columns, and their iterated Lie brackets
span $TX$ at every point.
\end{enumerate}
\end{theorem}

\subsection{Weighted energy}

Testing the weighted cofactor against $w^2$ produces a nonnegative gradient integral. Only distributional $\ddc$-closedness is needed: for a continuous $h>0$, the condition is
\begin{equation}\label{uni-eq:weightedclosed}
 \int_X\ddc\varphi\wedge hT=0\qquad(\varphi\in C^\infty(X,\mathbb R)).
\end{equation}

\begin{proof}[of Theorem~\ref{uni-thm:cofactor-intro}(1)]
The product rule gives
\begin{equation}\label{uni-eq:w2}
 \ddc(w^2)=2w\ddc w+2\mathrm{d}w\wedge \dc w
          =2w\ddc w+\frac{\mathrm{i}}{\pi}\partial w\wedge\bar\partial w.
\end{equation}
Approximate $w^2$ in $C^2$ by smooth functions in
\eqref{uni-eq:weightedclosed}.  Since $hT$ is continuous and
$\ddc w\wedge T=0$, this gives
\begin{equation}\label{uni-eq:energy}
 0=\frac1\pi\int_X h\,\mathrm{i}\partial w\wedge\bar\partial w\wedge T.
\end{equation}
The continuous integrand is nonnegative, so it vanishes pointwise.

For the nonnegative quadratic form $Q$,
$Q(\mathrm{d}w)=0$ implies $Q(\mathrm{d}w,\eta)=0$ for every $\eta$.  Therefore
\begin{equation}\label{uni-eq:annihilate}
 \mathrm{d}w_x\bigl(Q_x^\sharp\eta\bigr)=0\qquad(x\in X,\ \eta\in T_x^*X).
\end{equation}
The chain rule and \eqref{uni-eq:annihilate} make $w$ constant along every absolutely continuous curve whose velocity lies in $\operatorname{im}Q^\sharp$ almost everywhere.  One dense reachable set suffices by continuity.  If
instead $Q$ is positive definite on a dense set, \eqref{uni-eq:energy}
gives $\mathrm{d}w=0$ there and hence everywhere.
\end{proof}

The integration uses only the test function $w^2$. For general tests, integration by parts introduces derivatives of $hT$; $\ddc(hT)=0$ alone does not make $L\varphi=(\ddc\varphi\wedge T)/\mathrm{d}V$ a symmetric divergence operator.

\subsection{The real symbol and exterior-ball propagation}

To propagate an extremum without the weight, we write the linearized operator in real coordinates. Its symbol is continuous and semipositive. An exterior tangent ball to the maximum set can have its normal only in the kernel of that symbol; this forces the maximum set to be invariant under the range fields.

In a holomorphic chart $z_j=x_j+\mathrm{i} y_j$, let $C_{\rm cf}$ be the
Hermitian matrix characterized by
\begin{equation}\label{uni-eq:Ccf}
 \frac{\bigl(\mathrm{i}\sum H_{j\bar k}\mathrm{d}z_j\wedge \mathrm{d}\bar z_k\bigr)\wedge T}{\mathrm{d}V}
     =\operatorname{Tr}(C_{\rm cf}H)
\end{equation}
for every Hermitian matrix $H$.  Positivity of mixed wedges makes
$C_{\rm cf}$ positive semidefinite.  Write
$C_{\rm cf}=P_{\rm cf}+\mathrm{i} Q_{\rm cf}$, where $P_{\rm cf}$ is real
symmetric and $Q_{\rm cf}$ is real skew-symmetric.  In the real coordinate
order $(x_1,\ldots,x_n,y_1,\ldots,y_n)$,
\begin{equation}\label{uni-eq:real-symbol}
 L\varphi:=\frac{\ddc\varphi\wedge T}{\mathrm{d}V}
       =\operatorname{Tr}(aD^2\varphi),\qquad
 a=\frac1{8\pi}
 \begin{pmatrix}P_{\rm cf}&Q_{\rm cf}\\-Q_{\rm cf}&P_{\rm cf}\end{pmatrix}.
\end{equation}
\Needspace{8\baselineskip}
Indeed
$\varphi_{j\bar k}=(\varphi_{x_jx_k}+\varphi_{y_jy_k}
+\mathrm{i}\varphi_{x_jy_k}-\mathrm{i}\varphi_{y_jx_k})/4$; inserting this expression
in \eqref{uni-eq:Ccf} proves \eqref{uni-eq:real-symbol}.  For real coordinate vectors $\xi,\eta\in\mathbb R^n$, moreover,
\[
 (\xi,\eta)^Ta(\xi,\eta)
    =\frac1{8\pi}(\xi-\mathrm{i}\eta)^*C_{\rm cf}(\xi-\mathrm{i}\eta)\geq0.
\]
Thus $a$ is a continuous principal symbol, with no first-order term in
holomorphic coordinates, and its range is that of $Q^\sharp$.

An exterior-ball argument propagates the maximum for this operator.  Let $W\in C^2$
satisfy $LW\geq0$ in a chart, suppose $W\leq M$, and put
$K=\{W=M\}$.  If $z\in K$ is a nearest point to $y\notin K$, then
\begin{equation}\label{uni-eq:normalzero}
 (y-z)^Ta(z)(y-z)=0.
\end{equation}
To prove this, suppose the left side is positive, set
$r_0=|y-z|$, and shrink the tangent ball toward $z$:
$y'=z+\theta(y-z)$, $R=\theta r_0$, with $0<\theta<1$ small.
The closed ball $\overline{B(y',R)}$ lies in the chart and meets $K$
only at $z$.  This follows from the triangle inequality; equality in
$|k-y|\leq|k-y'|+|y'-y|\leq r_0$, for $k\in K$, forces $k=z$.
For large $\gamma$, the function
\[
 \psi(x)=\mathrm{e}^{-\gamma R^2}-\mathrm{e}^{-\gamma|x-y'|^2}
\]
satisfies
\[
 L\psi=2\gamma \mathrm{e}^{-\gamma|x-y'|^2}
   \bigl(\operatorname{tr}a-2\gamma(x-y')^Ta(x)(x-y')\bigr)<0
\]
in a neighborhood of $z$.  Choose a sufficiently small $r<R$, and
consider the lens $D=B(z,r)\cap B(y',R)$.  On its boundary lying on the
outer tangent sphere, $\psi=0$ and $W-M\leq0$.  On the remaining
compact boundary portion, $W-M<0$, since the tangent ball meets $K$
only at $z$.  Hence, for sufficiently small $\varepsilon>0$,
$W-M\leq\varepsilon\psi$ on $\partial D$.  A positive interior
maximum of $W-M-\varepsilon\psi$ would have nonpositive Hessian but
strictly positive image under $L$, which is impossible.  The same
inequality therefore holds in $D$.  Outside the tangent ball
$\psi\geq0$, so it also holds in a full neighborhood of $z$.
Second-order contact at $z$ now yields
\[
 0\leq LW(z)\leq\varepsilon L\psi(z)<0,
\]
proving \eqref{uni-eq:normalzero}.

Since $a\geq0$, its range is the orthogonal complement of its kernel.
Equation \eqref{uni-eq:normalzero} consequently implies
$(y-z)\cdot Z(z)=0$ for every range field $Z$.  If $Z$ is locally
Lipschitz and $\gamma$ is one of its flow curves starting in $K$, let
$d(t)$ denote its squared distance to $K$ in a sufficiently small
compact chart neighborhood.  Shortening the time interval ensures that
nearest points $z_t$ lie in the chart interior.  Write $D^+d(t)=\limsup_{h\downarrow0}(d(t+h)-d(t))/h$ for the upper right derivative. Holding such a nearest
point fixed when differentiating the distance from above gives
\begin{align*}
 D^+d(t)&\leq2(\gamma(t)-z_t)\cdot Z(\gamma(t))\\
 &=2(\gamma(t)-z_t)\cdot\bigl(Z(\gamma(t))-Z(z_t)\bigr)
 \leq2\operatorname{Lip}(Z)d(t).
\end{align*}
As $d(0)=0$, Gronwall's inequality gives $d(t)=0$.  The same argument
for $-Z$, followed by continuation, proves invariance under both signs
of time and under finite concatenations of such flows.

\Needspace{7\baselineskip}
\begin{proof}[of Theorem~\ref{uni-thm:cofactor-intro}(2)]
Apply the preceding argument to $W=w$ and to $W=-w$ in each chart.
Both are solutions of $LW=0$.  The maximum set of $w$ is nonempty,
closed, and contains every orbit starting in it.  If every orbit is dense,
this maximum set is all of $X$, proving constancy.  The argument for the
minimum set is identical.

For completeness, let a smooth family of range fields be bracket
generating.  At $p$, let $\mathcal D_p$ be the span of the values of
all original fields pushed forward by finite compositions of their local
flows and inverse flows.  This distribution contains the original fields
and is invariant under their flows.  Differentiation of pushed-forward
fields along a flow shows that it contains their nested Lie brackets;
finite-dimensional subspaces are closed, so the derivative remains in
$\mathcal D_p$.  Induction and the Jacobi identity give every iterated
bracket.  The assumed rank condition therefore gives
$\mathcal D_p=T_pX$.

Choose $2n$ such pushed-forward fields independent at $p$.  The flow
of each is a conjugate of an original flow and hence preserves its orbit.
The map obtained by successively flowing for $2n$ small times has
invertible derivative at zero.  The inverse function theorem puts a
neighborhood of $p$ in its orbit.  All orbits are open, and their
partition of connected $X$ has only one member.  This proves the stated
H\"ormander-type sufficient condition, without a constant-rank hypothesis
or a smooth matrix square root.  In particular, the columns of a smooth
symbol $a$ may be used as the original range fields.
\end{proof}

For $C^{2,1}$ potentials and a locally Lipschitz background, the cofactor columns are locally Lipschitz and the orbit criterion still applies. The bracket criterion uses smooth fields.

\Needspace{25\baselineskip}
\section{Gauduchon metrics and the cofactor}
\label{sec:gauduchon-surface-uniqueness}

Integration against $w^2$ also measures the obstruction to using the unweighted cofactor. For smooth potentials, the error is exactly the pairing of $w^2$ with $\ddc T$. On a surface the Gauduchon condition makes this error vanish for every pair; in dimension three an additional term depends on the potentials.

\begin{proposition}\label{prop:mixed-ddc-uniqueness}
Let $n\geq2$ and let $u,v$ be smooth $\omega$-plurisubharmonic functions with $A^n=B^n=f\,\mathrm{d}V$ for a smooth $f\geq0$.
Then
\begin{equation}\label{eq:mixed-uniqueness-defect}
 \int_X \mathrm{d}w\wedge\dc w\wedge T
       =\frac12\int_X w^2\ddc T.
\end{equation}
If $\ddc T=0$ and $\{f>0\}$ is dense in $X$, then $u-v$ is constant.
\end{proposition}
\begin{proof}
By \eqref{uni-eq:linearization}, $\ddc w\wedge T=0$. Integrating \eqref{uni-eq:w2} against $T$ and applying Stokes twice yields
\[
 \int_X\ddc(w^2)\wedge T=\int_Xw^2\ddc T.
\]
Substitution proves \eqref{eq:mixed-uniqueness-defect}.

For the uniqueness conclusion, apply \cref{uni-thm:cofactor-intro}(1) with $h=1$. At every point with $f>0$, the form $B$ is positive definite, so the term $B^{n-1}$ makes the cofactor strictly positive. Its positivity locus is therefore dense, as required.
\end{proof}

\Needspace{12\baselineskip}
\begin{corollary}\label{cor:gauduchon-surface-uniqueness}
Let $X$ be a connected compact complex surface and let $\omega$ be a Gauduchon metric. If $f\geq0$ is smooth and its zero set has empty interior, then two smooth $\omega$-plurisubharmonic solutions, both normalized to have supremum zero, of
\[
 (\omega+\ddc u)^2=cf\,\mathrm{d}V
\]
with the same $c>0$ coincide.
\end{corollary}
\begin{proof}
In dimension two, $T=A+B$ and
\[
 \ddc T=2\ddc\omega=0,
\]
which is exactly the Gauduchon condition on a surface. Smooth nonnegative $f$ with zero set of empty interior is positive on a dense set. The proposition gives equality up to a constant; normalization removes that constant.
\end{proof}

The next section removes the Gauduchon condition on a surface by integrating a signed quadratic difference. A conformal change of metric would change the equation and therefore cannot provide that extension.

In dimension three, direct expansion gives
\[
 \ddc T=3\ddc(\omega^2)+3\ddc\omega\wedge\ddc(u+v).
\]
The Gauduchon condition cancels the first term but need not cancel the second. For a concrete check, on $X=(\mathbb C/(\mathbb Z+\mathrm{i}\mathbb Z))^3$, with coordinates $z_1,z_2,z_3$ on the universal cover, write $e_j=\mathrm{i}\,\mathrm{d}z_j\wedge \mathrm{d}\bar z_j$, $x_j=\operatorname{Re}z_j$, and set
\[
 H=2+\cos(2\pi x_1),\qquad
 \omega=e_1+H e_2+H^{-1}e_3.
\]
The expansion
$\omega^2=2H e_1\wedge e_2+2H^{-1}e_1\wedge e_3+2e_2\wedge e_3$
has $\ddc\omega^2=0$: differentiating either variable coefficient adds a repeated $e_1$. Take $u=v=\delta\cos(2\pi x_2)$ with nonzero $\delta$ sufficiently small. Then $A=B>0$, but
\[
 \ddc T=6\ddc\omega\wedge\ddc u\ne0,
\]
because the $e_1\wedge e_3$ coefficient of $\ddc\omega$ and the $e_2$ coefficient of $\ddc u$ are both nonzero at $x_1=x_2=0$. Thus a Gauduchon equation background in dimension three can have a cofactor with $\ddc T\ne0$, even for equal strictly positive solution forms.

\section{Integrated rigidity in dimensions two and three}
\label{uni-sec:lowdim}

The linearized equation is not the only way to use equality of
determinants.  In low dimension a quadratic difference has a definite sign,
and its integral is fixed by exactness or by a sign on the torsion.

\begin{theorem}\label{uni-thm:sobolev}
Let $n\in\{2,3\}$ and let $u,v\in W^{2,n}(X)$.  Assume
$A,B\geq0$ almost everywhere and
\[
  A^n=B^n=f\,\mathrm{d}V,\qquad f>0\quad\hbox{almost everywhere}.
\]
For $n=2$ impose no differential condition on $\omega$.  For
$n=3$, assume
\begin{equation}
  \ddc\omega\wedge\mathrm{i}\eta\wedge\bar\eta\leq0
  \label{uni-eq:threefold-sign}
\end{equation}
for every $(1,0)$-covector $\eta$.  Then
$\ddc(u-v)=0$ distributionally and $u-v$ is almost everywhere constant.
For upper semicontinuous $\omega$-plurisubharmonic representatives the
equality holds pointwise.
\end{theorem}

\begin{proof}
Put
\[
 \begin{gathered}
 w=u-v,\qquad s=(u+v)/2,\qquad D=\ddc w=A-B,\\
 C=(A+B)/2=\omega+\ddc s.
 \end{gathered}
\]
The coefficients of $C,D$ lie in $L^n$; their top products are
integrable by H\"older's inequality.  At almost every point the common
determinant is positive, so $A,B$ are positive definite.  Diagonalize
them simultaneously by congruence, writing
$A=\sum e_j$, $B=\sum\lambda_j e_j$, where
$e_j=\mathrm{i}\theta_j\wedge\bar\theta_j$, $\lambda_j>0$, and
$\prod\lambda_j=1$.

In dimension two,
\begin{equation}\label{uni-eq:surface-sign}
 D^2=A^2+B^2-2A\wedge B
    =2(2-\lambda_1-\lambda_2)e_1\wedge e_2\leq0,
\end{equation}
with equality precisely when $A=B$.  In dimension three,
\begin{equation}\label{uni-eq:threefold-algebra}
 2C\wedge D^2=A^3+B^3-A^2\wedge B-A\wedge B^2\leq0.
\end{equation}
Indeed, with $E=e_1\wedge e_2\wedge e_3$, the four terms on the right
are
\[
 6E,\qquad 6E,\qquad 2\Bigl(\sum_j\lambda_j\Bigr)E,
 \qquad 2\Bigl(\sum_{i<j}\lambda_i\lambda_j\Bigr)E.
\]
Each sum is at least three, and equality again forces all
$\lambda_j=1$.

Choose smooth real $w_k\to w$ in $W^{2,n}$ and set
$D_k=\ddc w_k$.  Such approximations follow by a partition of unity,
local extension of compactly supported pieces, and convolution.  They
are not required to preserve semipositivity.  For $n=2$, Stokes gives
$\int_XD_k^2=0$, since
$D_k^2=\mathrm{d}(\dc w_k\wedge D_k)$.  Moreover
\[
 \|D_k^2-D^2\|_{L^1}
 \leq C\|D_k-D\|_{L^2}(\|D_k\|_{L^2}+\|D\|_{L^2})\longrightarrow0.
\]
Thus $\int_XD^2=0$, and the nonpositive integrand in
\eqref{uni-eq:surface-sign} vanishes almost everywhere.  Its strict equality
case gives $D=0$ almost everywhere.

For $n=3$, also choose smooth $s_k\to s$ in $W^{2,3}$, and put
$C_k=\omega+\ddc s_k$.  The needed identity is
\begin{equation}\label{uni-eq:torsion-integral}
 \int_X C_k\wedge D_k^2
 =\int_X\omega\wedge D_k^2
 =-\int_X\mathrm{d}w_k\wedge \dc w_k\wedge\ddc\omega.
\end{equation}
Here the first equality follows because
$\ddc s_k\wedge D_k^2=\mathrm{d}(\dc s_k\wedge D_k^2)$.
To check the sign in the second equality, put
$J_k=\mathrm{d}w_k\wedge \dc w_k$.  Using
$\dc \mathrm{d}=-\ddc $ and $(\dc )^2=0$,
\[
 \dc J_k=-D_k\wedge \dc w_k,\qquad \ddc J_k=-D_k^2.
\]
Stokes, first for $\mathrm{d}(\omega\wedge \dc J_k)$ and then for
$\dc (\mathrm{d}\omega\wedge J_k)$, gives
\begin{align*}
 \int_X\omega\wedge\ddc J_k&=-\int_X\mathrm{d}\omega\wedge \dc J_k,\\
 0&=-\int_X\ddc\omega\wedge J_k-\int_X\mathrm{d}\omega\wedge \dc J_k.
\end{align*}
Combining them yields \eqref{uni-eq:torsion-integral}.  Thus no unrecorded
first-order torsion term is discarded.

Strong $L^3$ convergence implies convergence of all three-factor
products in $L^1$.  For example,
\[
 \|(C_k-C)\wedge D_k^2\|_1
 \leq C\|C_k-C\|_3\|D_k\|_3^2\longrightarrow0,
\]
and
\[
 \|C\wedge(D_k^2-D^2)\|_1
 \leq C\|C\|_3\|D_k-D\|_3(\|D_k\|_3+\|D\|_3)\longrightarrow0.
\]
Also $\mathrm{d}w_k\wedge \dc w_k\to \mathrm{d}w\wedge \dc w$ in $L^{3/2}$, so the
right side of \eqref{uni-eq:torsion-integral} converges.  Consequently
\begin{equation}\label{uni-eq:torsion-limit}
 \int_XC\wedge D^2=-\int_X\mathrm{d}w\wedge \dc w\wedge\ddc\omega\geq0.
\end{equation}
The last sign uses
$\mathrm{d}w\wedge \dc w=(\mathrm{i}/(2\pi))\partial w\wedge\bar\partial w$
and \eqref{uni-eq:threefold-sign}.  The left integrand is nonpositive by
\eqref{uni-eq:threefold-algebra}, so it vanishes almost everywhere.  The
positive-density equality case gives $D=0$ almost everywhere.

In both dimensions $\ddc w=0$ as a distribution.  Taking the complex
trace in a chart gives the distributional Euclidean equation
$\Delta w=0$.  To recall why a smooth representative exists, convolve
locally with smooth radial kernels.  The convolutions are harmonic on
smaller balls and converge to $w$ in $L^1$.  Their mean-value property
against a fixed smaller radial kernel passes to the limit, identifying
$w$ almost everywhere with a smooth harmonic convolution there.  These
representatives agree on overlaps and give a global smooth harmonic
representative.  Its attained maximum on compact connected $X$ and the
local mean-value property make it constant.

If the original potentials are upper semicontinuous
$\omega$-plurisubharmonic representatives, each equals the limit of its
local ball averages: this follows from the submean inequality and upper
semicontinuity after adding a smooth local function.  The almost
everywhere constant-difference equality therefore holds pointwise.  Equal
mean normalizations give equality almost everywhere in the Sobolev class;
equal supremum normalizations give equality for the plurisubharmonic
representatives.
\end{proof}

\begin{corollary}\label{uni-cor:dense-lowdim}
For $C^2$ potentials in dimensions two and three, the corresponding
conclusion remains valid if the common continuous density is positive
merely on a dense set.  In dimension three retain
\eqref{uni-eq:threefold-sign}; in particular $\ddc\omega=0$ is sufficient.
\end{corollary}

\begin{proof}
At a point where the common top power vanishes, the algebraic expressions
are still nonpositive: in dimension two $D^2=-2A\wedge B$, and in
dimension three
$2C\wedge D^2=-A^2\wedge B-A\wedge B^2$.
The same integral identities force these continuous expressions to vanish
everywhere.  Their strict equality cases give $D=0$ on $\{f>0\}$.
Density and continuity then give $D=0$ globally.  No equality-case
assertion at a degenerate point is needed.
\end{proof}


The Hopf threefold of \cref{chap:hopf} illustrates the distinction between this uniqueness argument and a mass bound. Its metric satisfies $\ddc\omega=-\alpha^2$ with $\alpha\geq0$, so it meets the sign condition in \eqref{uni-eq:threefold-sign}. The Sobolev uniqueness theorem therefore applies on this example, even though its upper Monge--Amp\`ere volume is infinite.

\section{Bounded uniqueness and connectedness}
\label{sec:bounded-connected-uniqueness}

For a continuous density whose positive set is connected and dense, bounded solutions are unique up to a constant. An energy measure of $w=u-v$ replaces the pointwise gradient calculation available for smooth solutions. Its pushforward under $w$ has a Lipschitz density vanishing at the extreme values. If $w$ were nonconstant, these properties would force both extremal sets to have zero mass for the common Monge--Amp\`ere measure. A smooth test potential with a larger density near one of these sets then contradicts modified comparison.

\begin{theorem}\label{thm:connected-positive-uniqueness}\index{Uniqueness!connected positivity set}
Let $F\in C^0(X)$ satisfy $F\geq0$, $F\not\equiv0$, with $U=\{F>0\}$ connected and dense. If bounded $\omega$-plurisubharmonic functions $u,v$ satisfy
\[
\omega_u^n=\omega_v^n=F\omega^n,
\]
then $u-v$ is constant. In particular, $\sup_Xu=\sup_Xv$ implies $u=v$.
\end{theorem}

Continuity of $F$ supplies the local positive lower bounds and the neighborhoods of the extremal set used below. The auxiliary test densities are smooth even though $F$ is only continuous. Bounded solutions are continuous by \cref{cor:bounded-lp-continuity}, so we may use pointwise extrema throughout the proof. We shall also use modified comparison in the following form. If $b$ is continuous and $\omega_b\geq\varepsilon\omega$ for some $0<\varepsilon<1$, then $b/(1-\varepsilon)$ is $\omega$-plurisubharmonic, so \cref{thm:modified-comparison} applies with $(1-\varepsilon)v=b$. Thus its thin-sublevel estimate applies directly to such a strictly positive comparison potential $b$, and the correction factor tends to one as the height of the sublevel set tends to zero, with $\varepsilon$ fixed.

\begin{proposition}\label{prop:weak-extremum-localization}
Let $u,v$ be continuous $\omega$-plurisubharmonic functions with $\omega_u^n=\omega_v^n=F\omega^n$, where $F\geq0$ is bounded and measurable. Let $q$ be a smooth strictly $\omega$-plurisubharmonic function, and write $\omega_q^n=G\omega^n$. Then for every $0<\kappa<1$ and $\eta>0$,
\begin{equation}\label{eq:weak-localization-min}
F\omega^n\bigl(\{u-v<\min_X(u-v)+\eta\}\cap\{F>\kappa G\}\bigr)>0.
\end{equation}
The analogous assertion holds for the set $\{u-v>\max_X(u-v)-\eta\}$.
\end{proposition}

\begin{proof}
Choose $\delta\in(0,1]$ with $\omega_q\geq\delta \omega$. For $0<t<1$, put $b_t=(1-t)v+tq$, $m_t=\min_X(u-b_t)$, and $E_s=\{u<b_t+m_t+s\}$. Then $\omega_{b_t}\geq t\delta \omega$. Write $\mu=F\omega^n$, $d=\kappa^{-1/n}-1$, and $\gamma=nd$. By \eqref{eq:interior-minkowski-background},
\[
\omega_{b_t}^n\geq\bigl((1-t)F^{1/n}+tG^{1/n}\bigr)^n \omega^n.
\]
\Needspace{9\baselineskip}
On $L=\{F\leq\kappa G\}$ this is at least $(1+\gamma t)\mu$, since $G^{1/n}\geq(1+d)F^{1/n}$ there, and on $R=\{F>\kappa G\}$ it is at least $(1-nt)\mu$. Choose $s>0$ so small that the reverse comparison factor in \eqref{eq:modified-comparison}, with $\varepsilon=t\delta$, is at most $1+\gamma t/2$. Modified comparison on $E_s$, with $x=\mu(E_s\cap L)$ and $y=\mu(E_s\cap R)$, gives
\[
(1+\gamma t)x+(1-nt)y\leq(1+\gamma t/2)(x+y),
\qquad\text{and hence}\qquad y\geq\frac{1-\kappa^{1/n}}2\,\mu(E_s).
\]
The open set $E_s$ is nonempty by the definition of $m_t$. As $\omega_{b_t}^n\geq(t\delta)^n\omega^n$, \eqref{eq:modified-comparison} forces $\mu(E_s)>0$, hence $y>0$.

Let $a=\min_X(u-v)$ and $\Lambda=\|v-q\|_\infty$. Since $u-b_t=u-v+t(v-q)$, we have $|m_t-a|\leq t\Lambda$ and $E_s\subset\{u-v<a+2t\Lambda+s\}$. Given $\eta$, choose $t$ and then $s$ with $2t\Lambda+s<\eta$. Since $y=\mu(E_s\cap R)>0$, this proves \eqref{eq:weak-localization-min}. Interchanging $u$ and $v$ gives the assertion for the near-maximum set.
\end{proof}

\subsection{The energy of the difference}

In this subsection, $u,v$ are continuous $\omega$-plurisubharmonic functions with $A^n=B^n=F\omega^n$ for a bounded measurable function $F\geq0$, and $w=u-v$. Although the cofactor $T$ is now only a current, its product with $\mathrm dw\wedge\dc w$ still defines a positive measure. The identity $\ddc w\wedge T=0$ determines how this energy is distributed over the level sets of $w$.

\Needspace{24\baselineskip}
\begin{proposition}\label{prop:weak-level-energy}
There is a finite positive measure $E=\mathrm{d} w\wedge\dc w\wedge T$, defined below as a limit of regularizations, with the following properties.

For every smooth real function $\chi$ on $\mathbb R$,
\begin{equation}\label{eq:weak-energy-chain}
\int_X\chi''(w)\,\mathrm{d} E=\int_X\chi(w)\,\ddc T.
\end{equation}
The signed current $\ddc T$ is a finite measure. If $a=\min_Xw$, $b=\max_Xw$, and
\[
k(t)=\int_X(w-t)_+\,\ddc T,
\]
then $k$ is nonnegative, globally Lipschitz, supported on $[a,b]$, and
\begin{equation}\label{eq:weak-energy-pushforward}
w_*E=k(t)\,\mathrm{d} t.
\end{equation}
In particular, $E$ charges no level set of $w$.

Let $V$ be a relatively compact coordinate domain such that $F$ has a positive essential lower bound on a neighborhood of $\overline V$. Then there is $C_V<\infty$ such that every piecewise $C^1$ Lipschitz function $p$ on $[a,b]$ satisfies
\begin{equation}\label{eq:weak-local-coercivity}
\int_V|\mathrm{d}(p(w))|\,\mathrm{d}\lambda\leq C_V\left(\int_V|p'(w)|^2\,\mathrm{d} E\right)^{1/2},
\end{equation}
where $\mathrm{d}\lambda$ is Euclidean coordinate volume.
\end{proposition}

\begin{proof}
Locally, expand $T$ into powers of the smooth form $\omega$ multiplied by closed Bedford--Taylor products of $\ddc u$ and $\ddc v$. The operator $\ddc$ then acts only on the powers of $\omega$, and the local mass bounds in \cref{sec:psh-local-capacity} show that $\ddc T$ is a finite signed measure.

We first construct $E$ and prove \eqref{eq:weak-energy-chain}. On a coordinate ball, convolve $u$, $v$ and the coefficients of $\omega$ with a common nonnegative radial mollifier. Write the resulting smooth positive forms as $A_\rho=\omega_\rho+\ddc u_\rho$ and $B_\rho=\omega_\rho+\ddc v_\rho$, and let $T_\rho$ be their mixed sum. On a slightly larger ball, add the same fixed multiple of $|z|^2$ to each potential to obtain plurisubharmonic functions. Their radial regularizations decrease to the originals, so Bedford--Taylor continuity applies to all closed Hessian products. Since $u,v$ are continuous, $u_\rho,v_\rho,w_\rho$ converge uniformly on compact subsets.

For a closed smooth product $S_\rho$ appearing in $T_\rho$, the identity
\begin{equation}\label{eq:weak-energy-regularization}
\mathrm{d} w_\rho\wedge\dc w_\rho\wedge S_\rho
=\tfrac12\ddc(w_\rho^2S_\rho)-w_\rho\ddc w_\rho\wedge S_\rho
\end{equation}
has a distributional limit: both Hessian products on the right converge weakly with uniformly bounded local variation, and their weights converge uniformly. The same holds after multiplication by the convergent smooth factors coming from $\omega_\rho$. Thus the positive measures $E_\rho=\mathrm{d} w_\rho\wedge\dc w_\rho\wedge T_\rho$ converge to a finite positive measure $E$. By \eqref{eq:weak-energy-regularization}, this limit does not depend on the regularization, so the local limits agree on overlapping charts.

For smooth regularizations, the chain rule gives
\[
\ddc\chi(w_\rho)\wedge T_\rho
=\chi''(w_\rho)E_\rho+\chi'(w_\rho)\ddc w_\rho\wedge T_\rho.
\]
The last current tends to zero, because $\ddc w_\rho\wedge T_\rho=A_\rho^n-B_\rho^n$ tends weakly to $A^n-B^n=0$ with uniformly bounded variation, and $\chi'(w_\rho)$ converges uniformly. Passing to the limit and integrating by parts on the compact manifold $X$ gives \eqref{eq:weak-energy-chain}. A partition of unity justifies this passage without a global regularization; the terms containing derivatives of the cutoff functions cancel after summation.

Let $\sigma=w_*(\ddc T)$ and $\nu=w_*E$. Equation \eqref{eq:weak-energy-chain} says that $\nu''=\sigma$ as distributions on $\mathbb R$. Stokes' theorem and $\ddc w\wedge T=0$ give
\[
\sigma(\mathbb R)=\int_X\ddc T=0,\qquad
\int_{\mathbb R}s\,\mathrm{d}\sigma(s)=\int_Xw\,\ddc T=\int_X\ddc w\wedge T=0.
\]
Consequently $k(t)=\int(s-t)_+\,\mathrm{d}\sigma(s)$ is globally Lipschitz and vanishes outside $[a,b]$. It satisfies $k''=\sigma$, so the compactly supported distribution $\nu-k(t)\,\mathrm{d} t$ has second derivative zero and therefore vanishes. Positivity of $\nu$ and continuity of $k$ imply $k\geq0$. This proves \eqref{eq:weak-energy-pushforward}, and the assertion about level sets follows because $\nu$ is absolutely continuous.

We now prove \eqref{eq:weak-local-coercivity}. On a coordinate neighborhood of $\overline V$, set $\beta=\ddc|z|^2$ and choose $c_V>0$ with $F\omega^n\geq c_V\beta^n$. Let $A_{\mathrm{ac}}$ denote the absolutely continuous part of the current $A$, whose coefficients are locally integrable functions. Applying \cref{lem:interior-convolution} with the constant function $g=c_V^{1/n}$ and differentiating the coefficient measures, we obtain $\det_\beta A_{\mathrm{ac}}\geq c_V$ almost everywhere. Bounded plurisubharmonic functions lie in $W^{1,2}_{\mathrm{loc}}$ (\cref{sec:psh-local-capacity}), hence so does $w$ after the quadratic correction above, and its mollifications converge to it strongly in $W^{1,2}_{\mathrm{loc}}$, because the gradient of a mollification is the mollification of the gradient. Thus $w_\rho\to w$ strongly in $W^{1,2}_{\mathrm{loc}}$, and a subsequence of the gradients converges almost everywhere. Moreover, $A_\rho\to A_{\mathrm{ac}}$ almost everywhere, by differentiation of the coefficient measures.

\Needspace{6\baselineskip}
Since $T_\rho$ contains the term $A_\rho^{n-1}$ and all its terms are positive, Fatou's lemma, applied to $E_\rho$, yields the measure inequality
\begin{equation}\label{eq:weak-energy-ac-lower}
E\geq\mathrm{d} w\wedge\dc w\wedge A_{\mathrm{ac}}^{n-1}.
\end{equation}
\Needspace{10\baselineskip}
The cofactor matrix of $A_{\mathrm{ac}}$ has smallest eigenvalue at least $\det A_{\mathrm{ac}}/\operatorname{tr}A_{\mathrm{ac}}\geq c_V/\operatorname{tr}A_{\mathrm{ac}}$, where traces and determinants are taken with respect to $\beta$. The trace is locally integrable because $A$ has locally finite mass. The chain rule for Sobolev functions, \eqref{eq:weak-energy-ac-lower} and the Cauchy--Schwarz inequality give
\[
\int_V|\mathrm{d}(p(w))|\,\mathrm{d}\lambda
\leq C\left(\int_V|p'(w)|^2\,\mathrm{d} E\right)^{1/2}
\left(\int_V\frac{\operatorname{tr}A_{\mathrm{ac}}}{c_V}\,\mathrm{d}\lambda\right)^{1/2}.
\]
The last factor is finite and is absorbed into $C_V$. The values of $p'$ at the finitely many corners of $p$ are immaterial, because $E$ charges no level set of $w$.
\end{proof}

\Needspace{10\baselineskip}
\begin{lemma}\label{lem:weak-null-extremum-levels}
In the setting of \cref{prop:weak-level-energy}, assume in addition that $F$ is continuous and that $U=\{F>0\}$ is connected and dense. If $w=u-v$ is nonconstant, then
\[
F\omega^n(\{w=\min_Xw\})=F\omega^n(\{w=\max_Xw\})=0.
\]
\end{lemma}

\begin{proof}
Put $a=\min_Xw<b=\max_Xw$, and let $k$ be the function of \cref{prop:weak-level-energy}. It is Lipschitz and nonnegative, and $k(a)=0$. Choose $\epsilon_j\downarrow0$ with $a+\epsilon_j<b$. For $\delta>0$, set
\[
J_j(\delta)=\int_a^{a+\epsilon_j}\frac{\mathrm{d} t}{k(t)+\delta}.
\]
If $L$ is a Lipschitz constant for $k$, then $k(t)\leq L(t-a)$ on this interval, so $J_j(\delta)\to\infty$ as $\delta\downarrow0$. Choose $\delta_j>0$ with $J_j=J_j(\delta_j)\geq j$. Let $p_j=0$ on $(-\infty,a]$, $p_j=1$ on $[a+\epsilon_j,\infty)$, and
\[
p_j'(t)=\frac{1}{J_j(k(t)+\delta_j)}\quad(a<t<a+\epsilon_j).
\]
The pushforward identity \eqref{eq:weak-energy-pushforward} gives
\[
\int_X|p_j'(w)|^2\,\mathrm{d} E
=\frac{1}{J_j^2}\int_a^{a+\epsilon_j}\frac{k(t)}{(k(t)+\delta_j)^2}\,\mathrm{d} t
\leq\frac1{J_j}\longrightarrow0.
\]
Every point of $U$ has a coordinate neighborhood $V$ as in \cref{prop:weak-level-energy}, since $F$ is continuous. By \eqref{eq:weak-local-coercivity}, therefore, $\mathrm{d}(p_j(w))\to0$ in $L^1_{\mathrm{loc}}(U)$. The bounded functions $p_j(w)$ converge pointwise and in $L^1_{\mathrm{loc}}$ to $\mathbf1_{\{w>a\}}$, whose distributional gradient therefore vanishes on $U$. Since $U$ is connected, this indicator function is almost everywhere constant on $U$. As $w$ is nonconstant, $\{w>a\}$ is a nonempty open set, and its intersection with the dense open set $U$ is nonempty and open; hence the constant is one. Thus $\{w=a\}\cap U$ has zero volume, and since $F=0$ off $U$, the first assertion follows. Applying the argument to $-w$ gives the assertion for the maximum set.
\end{proof}

\Needspace{18\baselineskip}
\subsection{Concentrating smooth test potentials}

The next proposition produces smooth test potentials whose Monge--Amp\`ere density dominates a fixed multiple of a prescribed function $g$, with a constant independent of $\sup_Xg$ and of the derivatives of $g$. This uniformity allows the test densities to concentrate near a set of measure zero.

\Needspace{10\baselineskip}
\begin{proposition}\label{prop:weak-uniform-test-potentials}
For every $B_0>\|1\|_{L^2(\omega^n)}$ there is $\varepsilon=\varepsilon(X,\omega,B_0)>0$ such that every smooth function $g\geq1$ with $\|g\|_{L^2(\omega^n)}\leq B_0$ admits a smooth strictly $\omega$-plurisubharmonic function $q$ satisfying
\begin{equation}\label{eq:weak-uniform-test-determinant}
\omega_q^n\geq\varepsilon g \omega^n.
\end{equation}
\end{proposition}

\begin{proof}
Choose finitely many coordinate balls $B_j=\{|z_j|<R_j\}$, with closures in their charts, such that the concentric smaller balls $\{|z_j|<r_j\}$ cover $X$. On $B_j$, let $\psi_j$ solve the Dirichlet problem
\begin{equation}\label{eq:weak-local-dirichlet}
(\ddc\psi_j)^n=g\omega^n,\qquad \psi_j|_{\partial B_j}=0.
\end{equation}
The density of $g\omega^n$ with respect to coordinate volume is smooth and strictly positive. The classical smooth Dirichlet theorem therefore applies, since $C_j(|z_j|^2-R_j^2)$ is a smooth strict subsolution with zero boundary values for sufficiently large $C_j>0$. Thus $\psi_j$ is smooth and strictly plurisubharmonic on $\overline B_j$, and the maximum principle gives $\psi_j\leq0$.

We need a lower bound for $\psi_j$ independent of $\sup_Xg$. Write the coordinate equation as $\det(\partial_a\bar\partial_b\psi_j)=D_j$, and put $M_j=-\min_{B_j}\psi_j$. Let $x_j$ be a minimum point. For each real slope $p$ with $|p|<M_j/(2R_j)$, the function $\psi_j-p\cdot(x-x_j)$ attains its minimum at an interior point, because its value at $x_j$ is $-M_j$, while its boundary values are larger than $-M_j$. At such a contact point the real Hessian $S$ of $\psi_j$ is nonnegative, and the gradient of $\psi_j$ maps the contact set onto the ball of these slopes. The area formula and the matrix inequality
\[
\det_{\mathbb R}S\leq4^n\bigl(\det_{\mathbb C}(\partial_a\bar\partial_b\psi_j)\bigr)^2=4^nD_j^2
\]
therefore give $M_j\leq C(n,R_j)\|D_j\|_{L^2(B_j)}^{1/n}$. To prove the matrix inequality, let $J$ be the standard complex structure on the real coordinate space and average $S$ with $J^\mathsf TSJ$. Concavity of $\log\det$, and the fact that the eigenvalues of this $J$-invariant average occur in pairs, give the displayed factor; singular matrices are handled by approximation. Since the ratio of $\omega^n$ to coordinate volume is fixed and bounded, $\|D_j\|_2$ is controlled by $B_0$. Consequently
\begin{equation}\label{eq:weak-local-uniform-sup}
-M_0\leq\psi_j\leq0\quad\text{on every }B_j,
\end{equation}
where $M_0$ depends only on $X,\omega,B_0$ and the chosen finite cover.

Next we glue the local solutions. Put $\beta_j=\ddc|z_j|^2$. Choose $a_j>0$ such that $a_j\beta_j\leq \omega/2$ on $B_j$, and let $t_j=(R_j^2+r_j^2)/2$. Fix $\tau>0$ so small that $\tau M_0<a_j(R_j^2-r_j^2)/4$ for every $j$. The smooth function
\[
s_j=\tau\psi_j+a_j(t_j-|z_j|^2)
\]
is bounded below by a fixed positive constant on the smaller ball, negative on a collar of $\partial B_j$, and satisfies $\omega_{s_j}\geq \omega/2+\tau\ddc\psi_j$. Choose a smooth convex function $\phi_j\colon\mathbb R\to\mathbb R$ approximating $s\mapsto\max(s,0)$, equal to zero below a small negative threshold and to $s$ above the corresponding positive threshold, with $0\leq\phi_j'\leq1$ and $\phi_j''\geq0$; the thresholds are chosen smaller in absolute value than both margins of $s_j$. Set $Q_j=\phi_j(s_j)$ on $B_j$ and extend it by zero. This extension is smooth, and
\[
\omega_{Q_j}=\phi_j'(s_j)\omega_{s_j}+\bigl(1-\phi_j'(s_j)\bigr)\omega
+\phi_j''(s_j)\,\mathrm{d} s_j\wedge\dc s_j\geq \omega/2.
\]
\Needspace{8\baselineskip}
On the smaller ball, $Q_j=s_j$ and $\omega_{Q_j}\geq\tau\ddc\psi_j$. If $N$ is the number of balls, put $q=N^{-1}\sum_jQ_j$. Then $\omega_q=N^{-1}\sum_j\omega_{Q_j}\geq \omega/2$ everywhere. Each point belongs to at least one smaller ball $\{|z_j|<r_j\}$, so positivity of the other summands and \eqref{eq:weak-local-dirichlet} give
\[
\omega_q^n\geq(\tau/N)^n(\ddc\psi_j)^n=(\tau/N)^n g \omega^n.
\]
Thus \eqref{eq:weak-uniform-test-determinant} holds with $\varepsilon=(\tau/N)^n$, independent of $\sup g$ and of all derivatives of $g$.
\end{proof}

\begin{proof}[of \cref{thm:connected-positive-uniqueness}]
By \cref{cor:bounded-lp-continuity}, $u,v$ are continuous. In dimension one, $\ddc w=0$, so the maximum principle makes $w$ constant. Assume $n\geq2$.

Suppose that $w=u-v$ is nonconstant. Let $a=\min_Xw$, $K=\{w=a\}$, and $\mu=F\omega^n$. By \cref{lem:weak-null-extremum-levels}, $\mu(K)=0$. Choose $B_0>\|1\|_{L^2(\omega^n)}$ and let $\varepsilon>0$ be the constant from \cref{prop:weak-uniform-test-potentials}. Since $\mu(K)=0$, the compact set $K_0=K\cap\{F\geq\varepsilon/2\}$ has zero $\omega^n$-volume. Fix $M>1$ with $\varepsilon M>\max_XF$. By outer regularity of the smooth volume measure, choose an open neighborhood $O$ of $K_0$ so small that
\[
\int_X\omega^n+(M^2-1)\int_O\omega^n<B_0^2.
\]
Take a smooth $0\leq\chi\leq1$ equal to one near $K_0$ and supported in $O$, and set $g=1+(M-1)\chi$. If $K_0=\varnothing$, take $\chi=0$. Then $g\geq1$, $g=M$ on $K_0$, and $\|g\|_{L^2(\omega^n)}<B_0$, since $g^2\leq1+(M^2-1)\chi$.

By \cref{prop:weak-uniform-test-potentials}, there is a smooth strictly $\omega$-plurisubharmonic $q$ with $\omega_q^n=G\omega^n\geq\varepsilon g\omega^n$. On $K_0$ we have $G\geq\varepsilon M>F$; on $K\setminus K_0$ we have $F<\varepsilon/2<\varepsilon\leq G$. Thus $G>F$ throughout the compact set $K$. By compactness and continuity, we can choose $0<\kappa<1$ such that $F<\kappa G$ on $K$; put $W=\{F<\kappa G\}$. The complement of this open neighborhood of $K$ is compact, and $w-a$ has a strictly positive minimum there if it is nonempty. Hence some $\eta>0$ satisfies $\{w<a+\eta\}\subset W$.

\Cref{prop:weak-extremum-localization} asserts that $\mu(\{w<a+\eta\}\cap\{F>\kappa G\})>0$. By the preceding inclusion, however, this set is empty. This contradiction shows that $w$ is constant. If $\sup_Xu=\sup_Xv$, the constant is zero.
\end{proof}

\Needspace{16\baselineskip}
For a bounded measurable density, a positive lower bound gives the same conclusion. In this case coercivity holds on all of $X$, and the concentrating test potential can dominate the essential upper bound of the density on a whole neighborhood of the extremal set.

\begin{corollary}\label{cor:bounded-positive-density-uniqueness}
If $0<c_0\leq F\leq C_0<\infty$ almost everywhere and bounded $\omega$-plurisubharmonic functions $u,v$ satisfy $\omega_u^n=\omega_v^n=F\omega^n$, then $u-v$ is constant.
\end{corollary}
\begin{proof}
The functions $u,v$ are continuous by \cref{cor:bounded-lp-continuity}, and in dimension one the assertion follows from $\ddc(u-v)=0$ and the maximum principle. Let $n\geq2$, suppose that $w=u-v$ is nonconstant, and put $K=\{w=\min_Xw\}$. The proof of \cref{lem:weak-null-extremum-levels} applies with $U=X$, since \cref{prop:weak-level-energy} gives the pushforward identity and local coercivity on every chart. Thus $K$ has zero volume.

Fix $B_0>\|1\|_{L^2(\omega^n)}$ and the corresponding $\varepsilon$ in \cref{prop:weak-uniform-test-potentials}. Choose $M>1$ with $\varepsilon M>C_0$, and a smooth cutoff function $\chi$, equal to one on a neighborhood $W$ of $K$ and supported in an open set of volume so small that $g=1+(M-1)\chi$ satisfies $\|g\|_2<B_0$. The resulting test potential $q$ has density $G\geq\varepsilon g$. Choose $\kappa<1$ with $C_0<\kappa\varepsilon M$. Then $F<\kappa G$ almost everywhere on $W$. By continuity of $w$, a sufficiently thin set $\{w<\min_Xw+\eta\}$ lies inside $W$, which contradicts \cref{prop:weak-extremum-localization}. Hence $w$ is constant.
\end{proof}

\section{Bounded comparison on the equation background}
\label{uni-sec:bounded}

A condition on the background gives bounded uniqueness even when the positive-density region is disconnected. Ordinary comparison and an auxiliary solution whose density is enlarged on a strict sublevel set supply this criterion. The Guan--Li condition supplies ordinary comparison by making the total Monge--Amp\`ere mass independent of the potential.

Set $V=\int_X\omega^n$. We use the bounded products of \cref{prop:bounded-products} and the notation $\mathrm{MA}_\omega$.

The Guan--Li condition
\begin{equation}
   \ddc\omega=0,\qquad \ddc(\omega^2)=0
   \label{uni-eq:GL}
\end{equation}
will be imposed on the equation background $\omega$.

\Needspace{17\baselineskip}
\begin{theorem}\label{uni-thm:GL-intro}
Assume \eqref{uni-eq:GL}.  If $p>1$ and $f\geq0$ belongs to
$L^p(X,\omega^n)$ with positive integral, there is exactly one normalized
pair $(c,u)$, with $c>0$, $u\in\PSH(X,\omega)\cap L^\infty$,
$\sup_Xu=0$, satisfying
\[
  (\omega+\ddc u)^n=c f\,\omega^n.
\]
Moreover
\[
   c=\frac{\int_X\omega^n}{\int f\,\omega^n},
\]
and $u$ is continuous.  Thus two bounded solutions with the same
$L^p$ density differ by a constant; no positive lower bound on $f$ is
required.
\end{theorem}

We first establish the comparison principle through its equivalence with mass invariance.

\begin{theorem}\label{uni-thm:GL-equivalence}
The following conditions on the same equation form $\omega$ are
equivalent:
\begin{enumerate}
\item $\ddc\omega=0$ and $\ddc(\omega^2)=0$;
\item $\ddc\omega=0$ and $\mathrm{d}\omega\wedge \dc \omega=0$;
\item $\ddc(\omega^k)=0$ for $1\leq k\leq n-1$;
\item $\int_X\mathrm{MA}_\omega(a)=\int_X\omega^n$ for every bounded
      $\omega$-plurisubharmonic $a$;
\item for every bounded $\omega$-plurisubharmonic pair $a,b$,
\begin{equation}\label{uni-eq:ordinaryCP}
 \int_{\{a<b\}}\mathrm{MA}_\omega(b)
 \leq\int_{\{a<b\}}\mathrm{MA}_\omega(a).
\end{equation}
\end{enumerate}
The same equivalence holds when the tests in the last two conditions are
restricted to smooth potentials.  For $n=1$ all conditions are automatic;
for $n=2$ they reduce to $\ddc\omega=0$.
\end{theorem}

\begin{proof}
The differentiation formula
\begin{equation}\label{uni-eq:power-torsion}
 \ddc(\omega^k)=k\omega^{k-1}\wedge\ddc\omega
 +k(k-1)\omega^{k-2}\wedge \mathrm{d}\omega\wedge \dc \omega
\end{equation}
proves the equivalence of the first three conditions.  In dimensions one
and two the omitted equations vanish by degree, giving the stated
caveats.  The characterization of smooth mass invariance is the polarization argument of \cref{thm:mass-invariance-characterization}. Exact-background smooth decreasing regularization and bounded product continuity extend the identity to every bounded potential.

Here is a direct proof of bounded comparison from that identity. Set $q=\max(a,b)$. On $\{a<b\}$, locality gives $\mathrm{MA}_\omega(q)=\mathrm{MA}_\omega(b)$. On $\{a\geq b\}$, one has $a=q$ and $a\leq q$, so ordered contact (\cref{prop:ordered-contact-general}) gives
\[
 \mathbf1_{\{a\geq b\}}\mathrm{MA}_\omega(a)
 \leq \mathbf1_{\{a\geq b\}}\mathrm{MA}_\omega(q).
\]
All three total masses equal $V=\int_X\omega^n$. Subtracting the last inequality from this finite common mass proves \eqref{uni-eq:ordinaryCP}.

Conversely, bounded comparison implies bounded mass invariance: shift one of a bounded pair by a sufficiently large constant so
that its strict sublevel set is all of $X$, apply
\eqref{uni-eq:ordinaryCP}, and reverse the pair.  Comparing with the zero
potential yields the mass identity.  The bounded mass identity restricted
to smooth potentials implies the differential conditions by \cref{thm:mass-invariance-characterization}, closing all implications.
\end{proof}

For the uniqueness argument we need an auxiliary solution and the mixed inequality. Continuous existence supplies the first. If $G\geq0$ lies in $L^p(X,\omega^n)$ and $\int_XG\omega^n=V$, \cref{thm:kn-continuous} gives a continuous bounded solution with a positive normalizing constant. The mass identity makes that constant equal to one. Thus
\begin{equation}\label{uni-eq:aux-solvability}
 \mathrm{MA}_\omega(\rho)=G\omega^n
\end{equation}
is solvable, including for the indicator-truncated density below.

The bounded mixed inequality, proved from local Bedford--Taylor theory in \cref{lem:nonclosed-mixed-density}, applies to bounded $\omega$-plurisubharmonic functions $a,b$, nonnegative integrable densities $F,G$ and $0<t<1$. If $\mathrm{MA}_\omega(a)\geq F\omega^n$ and $\mathrm{MA}_\omega(b)\geq G\omega^n$, it gives
\begin{equation}\label{uni-eq:mixed-bounded}
 \mathrm{MA}_\omega((1-t)a+tb)
 \geq\bigl((1-t)F^{1/n}+tG^{1/n}\bigr)^n\omega^n,
 \qquad 0<t<1.
\end{equation}

\begin{proof}[of Theorem~\ref{uni-thm:GL-intro}]
The domination statement needed here is the case $c=0$ of \cref{thm:bounded-strict-domination}: if $\mathrm{MA}_\omega(a)(\{a<b\})=0$ for bounded $\omega$-plurisubharmonic functions $a,b$, then $a\geq b$.

Let $\mathrm{MA}_\omega(u)=\mathrm{MA}_\omega(v)=\mu=F\omega^n$, with
$F\in L^p$, $p>1$.  The mass identity gives $\mu(X)=V$.
If $u-v$ is nonconstant, shift $u$ so that its difference from $v$
has supremum $s>0$ and infimum $-s$.  Interchanging the pair if
necessary, arrange that $S=\{v<u\}$ satisfies $\mu(S)\leq V/2$.
Define
\[
 b_0=\frac{V-2\mu(S)}V\geq0,
 \qquad G=2F\mathbf1_S+b_0.
\]
This nonnegative $L^p$ function has total integral $V$, so choose the
continuous bounded solution $\rho$ of \eqref{uni-eq:aux-solvability}.
Take $C_0$ with $\rho-C_0\leq u$, and set
\[
 q_\delta=(1-\delta)u+\delta(\rho-C_0),\qquad
 E_\delta=\{v<q_\delta\}\subset S,
 \qquad 0<\delta<1.
\]
On $S$, the mixed inequality gives
\[
 \mathrm{MA}_\omega(q_\delta)\geq
 \bigl((1-\delta)+\delta2^{1/n}\bigr)^n\mu=:A_\delta\mu,
 \qquad A_\delta>1.
\]
Ordinary comparison on $E_\delta$ yields
\[
 A_\delta\mu(E_\delta)
 \leq\int_{E_\delta}\mathrm{MA}_\omega(q_\delta)
 \leq\int_{E_\delta}\mathrm{MA}_\omega(v)=\mu(E_\delta).
\]
Thus $\mu(E_\delta)=0$.  Bounded domination gives
$v\geq q_\delta$ everywhere; letting $\delta\downarrow0$ gives
$v\geq u$, contradicting $\sup_X(u-v)=s>0$.  Therefore the
difference is constant.  The argument never divides by $F$ and permits it to vanish on an open set.

Finally, for the density $f$ in the theorem put
$I=\int_Xf\omega^n>0$ and $c=V/I$.  Equation
\eqref{uni-eq:aux-solvability} supplies a continuous bounded solution for
$G=cf$, normalized by subtracting its supremum.  Any competing bounded
solution with constant $c'$ satisfies $V=c'I$, hence $c'=c$;
the just-proved uniqueness and the common normalization identify the
potentials.  In particular every bounded solution with this $L^p$
density is continuous, even though continuity was not assumed.
\end{proof}

\begin{remark}
Theorem~\ref{uni-thm:GL-equivalence} is an equivalence of differential,
comparison, and mass-invariance properties.  It is not an equivalence
between the Guan--Li condition and uniqueness.  The other criteria in this chapter apply without that background condition.  Conversely, imposing it
on an auxiliary metric rather than on $\omega$ does not justify
\eqref{uni-eq:ordinaryCP} for $\mathrm{MA}_\omega$.
\end{remark}

\Needspace{12\baselineskip}
\section{Smooth nonuniqueness for a degenerate density}
\label{sec:degenerate-hermitian-nonuniqueness}

The first example puts the freedom in the potential inside an open region where the density is zero. Three independent rank-one directions reduce equality of the Monge--Amp\`ere measures to equality of scalar products. We construct two semipositive forms with the same cube, then add an exact correction to obtain a strictly positive equation background.

Let $X=\mathbb P(\mathbb C^2\oplus\mathbb C^2)=\mathbb P^3$, with block homogeneous coordinates $[A:B]$, where $A,B\in\mathbb C^2$, and let $\omega_X$ be its Fubini--Study form.

\begin{theorem}\label{thm:degenerate-hermitian-nonuniqueness}
On $X=\mathbb P^3$ there exist a smooth Hermitian form $\omega$ with $\mathrm{d}\omega\ne0$, a smooth nonnegative nonzero function $f$, and distinct smooth functions $u,v\in\PSH(X,\omega)$ such that
\[
 \sup_Xu=\sup_Xv=0,\qquad
 (\omega+\ddc u)^3=f\omega_X^3=(\omega+\ddc v)^3.
\]
The density vanishes on a nonempty open set. Both potentials have minimal singularities.
\end{theorem}

\subsection{Three rank-one directions}
\label{subsec:numa-rank-directions}

On $U=\{A\ne0,B\ne0\}$ put
\[
 r=\log\frac{|A|^2}{|B|^2},\qquad
 P=\ddc\log|A|^2,\quad Q=\ddc\log|B|^2,\quad
 R=\mathrm{d}r\wedge \dc r.
\]
The logarithms individually use a local homogeneous lift; their Hessians are independent of that lift. The forms $P,Q$ are pullbacks of Fubini--Study forms from the two projective lines. All three forms are semipositive of rank one, and
\begin{equation}\label{eq:numa-rank-frame}
 P^2=Q^2=R^2=0,\qquad P\wedge Q\wedge R>0.
\end{equation}
Indeed, in local coordinates $A=(1,a)$ and $B=\tau(1,b)$, $P$ and $Q$ are positive multiples of $\mathrm{i}\,\mathrm{d}a\wedge \mathrm{d}\bar a$ and $\mathrm{i}\,\mathrm{d}b\wedge \mathrm{d}\bar b$, while
\[
 \partial r=\frac{\bar a\,\mathrm{d}a}{1+|a|^2}
             -\frac{\mathrm{d}\tau}{\tau}
             -\frac{\bar b\,\mathrm{d}b}{1+|b|^2}.
\]
\Needspace{6\baselineskip}
Its $\mathrm{d}\tau$ component is nonzero. The three rank-one directions are therefore independent. In particular, for nonnegative $x,y,z$,
\begin{equation}\label{eq:numa-cube}
 (xP+yQ+zR)^3=6xyz\,P\wedge Q\wedge R,
\end{equation}
and the form is positive exactly when all three coefficients are positive.

\Needspace{5\baselineskip}
For a smooth function $G:\mathbb R\to\mathbb R$ the chain rule is
\begin{equation}\label{eq:numa-chain}
 \ddc G(r)=G'(r)(P-Q)+G''(r)R.
\end{equation}
\Needspace{7\baselineskip}
Let $\sigma(r)=\mathrm{e}^r/(1+\mathrm{e}^r)$. Since
$\log(|A|^2+|B|^2)=\log|B|^2+\log(1+\mathrm{e}^r)$, we have
\begin{equation}\label{eq:numa-fs}
 \omega_X=\sigma(r)P+\sigma(-r)Q+\sigma'(r)R.
\end{equation}
A function of $r$ constant at each end of $\mathbb R$ extends smoothly over $X\setminus U$. A form whose coefficients equal those of \eqref{eq:numa-fs} for $|r|\geq2$ extends smoothly by $\omega_X$ there.

\subsection{The scalar coefficients}
\label{subsec:numa-scalar-coefficients}

The central coefficients will make a different rank-one factor vanish in each solution form. Outside the central interval they will return to the Fubini--Study coefficients, so the resulting forms extend across the omitted projective lines. Primes denote derivatives with respect to $r$.

Choose an even smooth bump $\beta$ supported on $[-1,1]$, strictly positive on $(-1,1)$ and with $\beta'>0$ on $(-1,0)$, normalized by $\int_{-1}^1\beta=2$. A positive multiple of $\exp(-1/(1-r^2))$ on $(-1,1)$, extended by zero, does this. Define
\[
 F(r)=-1+\int_{-1}^r\beta(t)\,\mathrm{d}t\quad(-1<r<1),
 \qquad F=-1\ (r\leq-1),\quad F=1\ (r\geq1),
\]
and put $p=\beta/2$, $q=\beta'/2$. Then $F$ is smooth and odd, $p$ is even, and $q$ is odd, positive on $(-1,0)$ and negative on $(0,1)$. Both $p$ and $q$ vanish to infinite order at the endpoints.

Take smooth $a\geq0$ with $a=1$ on $(-\infty,-2/3]$ and $a=0$ on $[-1/3,\infty)$, and put $b(r)=a(-r)$. Choose an even nonnegative $\chi\in C_c^\infty((-1/3,1/3))$ with $\chi(0)>0$. On $[-1,1]$ set
\[
 z_0=(q^2+\chi^2)^{1/2}.
\]
It is smooth: the expression under the square root is positive in the interior; near each endpoint $z_0=|q|$ with fixed sign of $q$. It is flat at the endpoints and positive in the interior. Also,
\begin{equation}\label{eq:numa-coefficient-cancellation}
 z_0\geq|q|,\qquad a(z_0-q)=0,\qquad b(z_0+q)=0.
\end{equation}
The cancellations follow from the disjoint supports of $a$ and $\chi$, and from the sign of $q$ on the support of $a$; the other equality follows by reflection.

\Needspace{12\baselineskip}
On $(-\infty,-1]$ choose smooth functions $\ell,s$ such that $0\leq\ell\leq1$, $\ell=0$ on $(-\infty,-2]$, $\ell=1$ on $[-3/2,-1]$, while $s>0$ for $r<-1$, $s=1$ on $(-\infty,-3/2]$, and $s$ is flat with value zero at $-1$. Such functions are obtained from ordinary smooth cutoffs and a flat positive endpoint function. Define
\[
 \begin{aligned}
 x&=(1-\ell)\sigma+\ell,&
 y&=(1-\ell)\sigma(-r)+\ell s,&
 z&=(1-\ell)\sigma'+\ell s &&(r\leq-1),\\
 x&=p+a,& y&=p+b,& z&=z_0 &&(-1\leq r\leq1),
 \end{aligned}
\]
and for $r\geq1$ use $x(r)=y(-r)$, $y(r)=x(-r)$, $z(r)=z(-r)$.

The definitions join smoothly. At $-1$ their values are $(1,0,0)$, with all discrepancies flat; near that endpoint $a=1$ and $b=0$. The central definitions are symmetric under reflection, so the same holds at $1$. We have
\[
 x,y,z>0\quad\hbox{for }r\ne\pm1,\qquad
 (x,y,z)(-1)=(1,0,0),\quad (x,y,z)(1)=(0,1,0).
\]
For $|r|\geq2$ their values are $(\sigma(r),\sigma(-r),\sigma'(r))$.

\subsection{Two forms with the same cube}
\label{subsec:numa-equal-cubes}

The cancellation identities now turn the scalar coefficients into two semipositive forms with the same cube. Their average will be strictly positive except at the two ends of the central interval.

The globally smooth form and function
\[
 \omega_0=x(r)P+y(r)Q+z(r)R,\qquad w=F(r)
\]
are understood by the extensions above. The form is semipositive, and positive except on $M_- =\{r=-1\}$ and $M_+=\{r=1\}$. Put
\[
 H=\tfrac12\ddc w=p(P-Q)+qR,\qquad A_0=\omega_0-H,\quad B_0=\omega_0+H.
\]
For $|r|\geq1$ we have $H=0$. On $|r|<1$,
\[
 A_0=aP+(2p+b)Q+(z_0-q)R,\qquad
 B_0=(2p+a)P+bQ+(z_0+q)R.
\]
All coefficients are nonnegative. The cancellations in \eqref{eq:numa-coefficient-cancellation} give $A_0^3=B_0^3=0$ on this band. Outside it both cubes equal $\omega_0^3$. Thus $A_0,B_0$ are globally smooth semipositive forms with equal cubes. That common top-degree form is smooth, vanishes on the closed band $|r|\leq1$, and is positive on its complement, including near $X\setminus U$.

\Needspace{12\baselineskip}
\subsection{Restoring strict positivity of the background}
\label{subsec:numa-positive-background}

The forms $A_0,B_0$ already have the required common measure. We obtain a Hermitian background by adding the same exact correction to their average and subtracting its potential from both solutions. The correction must restore positivity at $M_\pm$ while preserving it in the transition regions.

Take $\kappa\in C_c^\infty((-2,2))$, equal to one on $[-3/2,3/2]$. For a parameter $\varepsilon>0$ to be chosen, set
\[
 \rho(r)=\varepsilon(r^2-1)\kappa(r),\qquad
 \omega=\omega_0+\ddc\rho.
\]
Its coefficients on $U$ are $x+\rho'$, $y-\rho'$ and $z+\rho''$. On $[-3/2,3/2]$ one has $\rho'=2\varepsilon r$ and $\rho''=2\varepsilon$. The last coefficient is positive. For $r\geq0$, $x+2\varepsilon r$ is positive, including at $r=0$ since $x(0)=p(0)>0$. On $[-3/2,0]$, the function $x$ has a positive minimum $m$, so $x+2\varepsilon r\geq m-3\varepsilon>0$ for small $\varepsilon$. Reflection gives the corresponding estimate for $y-2\varepsilon r$. On the compact region $3/2\leq|r|\leq2$, all three original coefficients have a positive minimum and the perturbations are $O(\varepsilon)$. Outside that region and the central interval the form is exactly $\omega_X$. This proves strict positivity globally.

For completeness, nonclosedness can be checked directly. Since $\mathrm{d}R=-\mathrm{d}r\wedge(P-Q)$,
\begin{equation}\label{eq:numa-differential}
 \mathrm{d}\omega=\mathrm{d}\omega_0
   =\mathrm{d}r\wedge\bigl((x'-z)P+(y'+z)Q\bigr).
\end{equation}
At $r=0$, $x'=y'=0$ and $z=\chi(0)>0$. Hence $\mathrm{d}\omega=z\,\mathrm{d}r\wedge(Q-P)\ne0$ there. The nonzero $\mathrm{d}\tau$ component used in \eqref{eq:numa-rank-frame} also shows that $\mathrm{d}r\wedge P$ and $\mathrm{d}r\wedge Q$ are linearly independent.

\begin{proof}[of \cref{thm:degenerate-hermitian-nonuniqueness}]
Set $\widetilde u=-w/2-\rho$ and $\widetilde v=w/2-\rho$. The construction gives
\[
 \omega+\ddc\widetilde u=A_0\geq0,\qquad
 \omega+\ddc\widetilde v=B_0\geq0.
\]
Subtract their respective suprema to obtain $u,v$. Their Hessians are unchanged, and $v-u=w+\text{constant}$ is nonconstant. The function $f=A_0^3/\omega_X^3=B_0^3/\omega_X^3$ is smooth, nonnegative and nonzero. It vanishes on the nonempty open band $|r|<1$. Both smooth potentials have minimal singularities because $V_\omega=0$ for a Hermitian background.
\end{proof}

Every Hermitian background on $\mathbb P^3$ has bounded mass: the Fubini--Study form has invariant mass by closed bounded Stokes, and \cref{prop:bmp-metric} makes finiteness independent of the metric. This applies to all the projective backgrounds constructed below.

The linearized difference also explains the failure of the smooth maximum argument. The tensor
\[
 T=A_0^2+A_0\wedge B_0+B_0^2
\]
satisfies $\ddc w\wedge T=B_0^3-A_0^3=0$, but $T=0$ on $M_\pm$, where $A_0=B_0=\omega_0$ has rank one. The resulting equation is not globally elliptic, even though both potentials and their common density are smooth.

\section{A family of solutions from the radial equation}
\label{sec:numa-radial-family}

The open zero-density band gives more freedom than the two smooth solutions exhibit. Writing a further potential as $G(r)-\rho$ reduces the equation to choosing a slope $G'$ for which one of the three factors vanishes. Switching the vanishing factor at two suitable levels produces a continuum of $C^{1,1}$ solutions for the same background $\omega$ and density $f$.

Retain the coefficients and correction of the preceding construction. On $-1<r<1$ we have
\[
 x=p+a,\qquad y=p+b,\qquad z=z_0,
 \qquad p'=q,\quad z_0\geq|q|.
\]
Both $p$ and $z_0$ are even, $p>0$ in the interior, and $p'<0$ on $(0,1)$. The strict inequality $z_0>|q|$ holds on a neighborhood of zero. Hence
\begin{equation}\label{eq:numa-half-excess}
 \int_0^1z_0(r)\,\mathrm{d}r>\int_0^1|p'(r)|\,\mathrm{d}r=p(0).
\end{equation}

\begin{proposition}\label{prop:numa-radial-continuum}
The same smooth background and smooth density in \cref{thm:degenerate-hermitian-nonuniqueness} admit uncountably many distinct normalized $C^{1,1}$ solutions of
\[
 \mathrm{MA}_\omega(u)=f\omega_X^3.
\]
They include a directly constructed solution invariant under the block exchange $\iota[A:B]=[B:A]$, if the cutoff defining $\rho$ is chosen even.
\end{proposition}
\Needspace{6\baselineskip}
\begin{proof}
Choose $-1<I<0<J<1$ satisfying
\begin{equation}\label{eq:numa-transition-balance}
 \int_I^J z_0(r)\,\mathrm{d}r=p(I)+p(J).
\end{equation}
We first show there are an interval's worth of choices. For fixed $I$, define
\[
 K(I,J)=\int_I^Jz_0(r)\,\mathrm{d}r-p(I)-p(J).
\]
At $I=0$ we have $K(0,0)=-2p(0)<0$ and $K(0,1)>0$ by \eqref{eq:numa-half-excess}. Continuity gives $\epsilon>0$ such that these opposite signs persist for every $I\in[-\epsilon,0]$. For $0<J<1$,
\[
 \partial_JK(I,J)=z_0(J)-p'(J)>0.
\]
The intermediate value theorem therefore supplies a unique $J=J(I)\in(0,1)$ for every $I\in[-\epsilon,0)$. Shrink $\epsilon$ below one.

For any pair satisfying \eqref{eq:numa-transition-balance}, define the slope on $[-1,1]$ by
\[
 h_{I,J}(r)=
 \begin{cases}
 p(r),&-1\leq r\leq I,\\
 p(I)-\displaystyle\int_I^r z_0(s)\,\mathrm{d}s,&I\leq r\leq J,\\
 -p(r),&J\leq r\leq1,
 \end{cases}
\]
and extend it by zero outside that interval. The balance identity makes it continuous at $J$, and continuity at $I$ is immediate. Each piece has bounded derivative. Thus $h_{I,J}$ is globally Lipschitz; it is flat at $\pm1$. Let $G_{I,J}$ be any primitive. It is $C^{1,1}$ and constant at each end of the real line. Consequently $G_{I,J}(r)$ extends to a global $C^{1,1}$ function on $X$, smooth near $X\setminus U$.

On the middle interval the slope lies between $-p$ and $p$. Indeed,
\[
 (h_{I,J}-p)'=-z_0-p'\leq0,
 \qquad
 (h_{I,J}+p)'=-z_0+p'\leq0.
\]
The first difference starts at zero at $I$; the second ends at zero at $J$. These two inequalities give $h_{I,J}\leq p$ and $h_{I,J}\geq-p$, respectively.

Put $\widetilde u_{I,J}=G_{I,J}(r)-\rho$. Its solution form is, almost everywhere on $U$,
\begin{equation}\label{eq:numa-slope-form}
 \omega+\ddc\widetilde u_{I,J}
 =\omega_0+\ddc G_{I,J}(r)
 =(x+h_{I,J})P+(y-h_{I,J})Q+(z+h_{I,J}')R.
\end{equation}
On $I<r<J$, the first two coefficients are nonnegative because $x=p+a$, $y=p+b$ and $-p\leq h_{I,J}\leq p$. The third is zero. On $-1<r<I$, we have $h_{I,J}=p$ and $b=0$, since $I<0$; the second coefficient is zero and the others are nonnegative, using $z_0+p'\geq0$. On $J<r<1$, we have $h_{I,J}=-p$ and $a=0$; the first coefficient is zero and the others are nonnegative, using $z_0-p'\geq0$. Outside the band, $h_{I,J}=h_{I,J}'=0$, so the form equals $\omega_0$.

There is no additional hypersurface term in \eqref{eq:numa-slope-form}: $G_{I,J}$ is $C^{1,1}$, and its first derivative is continuous at both transition points. Its distributional Hessian is the bounded almost-everywhere Hessian. The preceding coefficient inequalities therefore prove $\omega+\ddc\widetilde u_{I,J}\geq0$ as a current.

By \cref{prop:c11-products}, the bounded product of this $C^{1,1}$ potential agrees with the almost-everywhere wedge of its Hessian.

The rank-one formula \eqref{eq:numa-rank-frame} now shows that the cube of \eqref{eq:numa-slope-form} is zero almost everywhere on the band, and equals $\omega_0^3$ off it. These are exactly the defining values of $f\omega_X^3$ in the preceding construction. Thus $\widetilde u_{I,J}$ solves the equation. Subtracting its supremum gives a normalized solution $u_{I,J}$.

Different $I$ in the interval above give different normalized solutions. If $I_1<I_2<0$ and $I_1<r<I_2$, then the slope belonging to $I_2$ is $p(r)$, whereas
\[
 p(r)-h_{I_1,J(I_1)}(r)
       =\int_{I_1}^r(p'(s)+z_0(s))\,\mathrm{d}s>0,
\]
because $p'>0$ on $(-1,0)$. The two primitives therefore cannot differ by a constant. An interval contains uncountably many parameters, proving the first assertion. None of these slopes is identically $p$ or identically $-p$, so the constructed solutions are also distinct from the original two smooth solutions.

Finally, the function
\[
 L(t)=\int_0^t z_0(r)\,\mathrm{d}r-p(t),\qquad 0\leq t\leq1,
\]
is negative at zero, positive at one, and strictly increasing on $(0,1)$. Its unique zero $t_*$ gives an admissible symmetric pair $I=-t_*$, $J=t_*$. Evenness of $z_0$ and $p$ then makes $h_{I,J}$ odd: on the middle interval it is $-\int_0^r z_0(s)\,\mathrm{d}s$. Choose its primitive even. If $\rho$ is even, $G_{I,J}(r)-\rho$ and its normalized version are fixed by $\iota$. This proves the explicit invariant-solution assertion.
\end{proof}

The transition hypersurfaces lie inside the zero-density band. The resulting family also raises a selection question for positive-exponent approximation. For the symmetric data, \cref{thm:numa-selected-limit} proves that it selects the symmetric member just constructed.

\section{Nonuniqueness with a zero set of empty interior}
\label{sec:numa-thin-zero-set}

We now keep the Hermitian background of \cref{thm:degenerate-hermitian-nonuniqueness}, but change the density so that it is positive off $M_-\cup M_+$, where $M_\pm=\{r=\pm1\}$. For potentials of the form $G(r)-\rho$, the equation on the middle band becomes a first-order ordinary differential equation for the slope $G'$. Barriers keep its solutions between the admissible slopes, and infinite-order vanishing at the endpoints then makes their extensions by zero smooth.

\Needspace{9\baselineskip}
\begin{theorem}\label{thm:numa-thin-zero-set}\index{Zero set!thin degeneracy set}
There are a smooth nonclosed Hermitian form $\omega$ on $\mathbb P^3$ and a smooth nonnegative density $g$, positive off two disjoint smooth real hypersurfaces, such that the equation
\[
 \mathrm{MA}_\omega(u)=g\omega_X^3
\]
has uncountably many distinct smooth solutions normalized by $\sup_Xu=0$. The background $\omega$ has bounded mass.
\end{theorem}

Thus a zero set with empty interior, even one of measure zero, does not ensure uniqueness of smooth normalized solutions.

\Needspace{12\baselineskip}
We use the coefficients of \cref{subsec:numa-scalar-coefficients}, with the explicit choice
\[
 p(r)=C\exp\!\left(-\frac1{1-r^2}\right),\qquad -1<r<1,
 \qquad \int_{-1}^1p(r)\,\mathrm{d}r=1,
\]
extended by zero, where the integral condition determines $C>0$; then $q=p'$. Fix $0<s<1$, and set
\[
 k_s(r)=(x(r)+sp(r))(y(r)-sp(r))(z(r)+sq(r)).
\]
For $|r|>1$ this is $xyz$. If $G'=sp$, then the potential $G(r)-\rho$ has $\omega+\ddc(G(r)-\rho)=\omega_0+sH$. This form is strictly positive for $-1<r<1$, since
\[
 x+sp=a+(1+s)p>0,\quad
 y-sp=b+(1-s)p>0,\quad
 z_0+sq\geq(1-s)z_0>0.
\]
It is also strictly positive outside the closed band, and at $r=\pm1$ it equals the rank-one form $\omega_0$. Consequently
\begin{equation}\label{eq:numa-thin-density}
 g\omega_X^3=(\omega_0+sH)^3
\end{equation}
defines a smooth density $g$ that is positive exactly off $M_-\cup M_+$. These hypersurfaces have measure zero and empty interior. On the band $|r|<1$, a potential $G(r)-\rho$ with slope $h=G'$ solves the equation exactly when $(x+h)(y-h)(z+h')=k_s$, which is the differential equation \eqref{eq:numa-thin-ode} below.

\Needspace{14\baselineskip}
\begin{lemma}\label{lem:numa-thin-trapping}
There exists $\eta_0>0$ such that for every $0<\eta<\eta_0$ the initial-value problem
\begin{equation}\label{eq:numa-thin-ode}
 h'=\frac{k_s(r)}{(x(r)+h)(y(r)-h)}-z(r),
 \qquad h(0)=sp(0)-\eta,
\end{equation}
has a smooth solution on $(-1,1)$ with $-x<h<y$ and $h<sp$. Moreover, on some fixed endpoint neighborhoods there are constants $0<c_L<s<t<1$, independent of $\eta$, such that
\[
 c_Lp<h<sp\quad\hbox{near }-1,
 \qquad -tp<h<sp\quad\hbox{near }1.
\]
\end{lemma}
\begin{proof}
The function $h_s=sp$ solves \eqref{eq:numa-thin-ode} and lies strictly between $-x$ and $y$ throughout $(-1,1)$. Write $F_s(r,h)$ for the right-hand side of the equation, which is smooth on the open domain $-x(r)<h<y(r)$.

Near $-1$ the coefficients are exactly
\[
 x=1+p,\qquad y=p,\qquad z=q>0.
\]
For a constant $c\in(-1,1)$, the product of coefficients for the trial slope $cp$ is
\[
 k_c=(1-c^2)\bigl(1+(1+c)p\bigr)pq.
\]
Fix $c_L\in(0,s)$. At $p=0$, the factor preceding $pq$ for $c=c_L$ exceeds that for $c=s$ by $s^2-c_L^2>0$. Thus, on a sufficiently small fixed neighborhood of the left endpoint, $k_{c_L}>k_s$ and
\begin{equation}\label{eq:numa-left-trap}
 F_s(r,c_Lp)-(c_Lp)'=
 \frac{k_s-k_{c_L}}{(x+c_Lp)(y-c_Lp)}<0.
\end{equation}
Thus $c_Lp$ is a lower barrier when the equation is solved backward, that is, toward $-1$.

Near $1$ the exact coefficients are
\[
 x=p,\qquad y=1+p,\qquad z=-q>0,
\]
and the trial slope $cp$ has product of coefficients
\[
 k_c=(1-c^2)\bigl(1+(1-c)p\bigr)p(-q).
\]
Fix $s<t<1$. At $p=0$, the factor for $c=s$ exceeds that for $c=-t$ by $t^2-s^2>0$. On a sufficiently small fixed neighborhood of the right endpoint, therefore, $k_s>k_{-t}$, so
\begin{equation}\label{eq:numa-right-trap}
 F_s(r,-tp)-(-tp)'=
 \frac{k_s-k_{-t}}{(x-tp)(y+tp)}>0.
\end{equation}
Thus $-tp$ is a lower barrier when the equation is solved forward, toward $1$.

Choose $r_L<0<r_R$ in these respective neighborhoods. The graph of $h_s$ over $[r_L,r_R]$ is a compact subset of the domain of $F_s$. Existence and continuous dependence on initial values for smooth ordinary differential equations therefore give $\eta_0>0$ such that, for every $0<\eta<\eta_0$, the solution with initial value $sp(0)-\eta$ exists on $[r_L,r_R]$, stays in the domain, and satisfies
\[
 h(r_L)>c_Lp(r_L),\qquad h(r_R)>-tp(r_R).
\]
By local uniqueness, two solutions of the equation cannot cross. Since $h(0)<h_s(0)$, we have $h<h_s$ wherever $h$ exists.

Follow the solution backwards from $r_L$; it cannot cross $c_Lp$. Indeed, at a first crossing point the difference $h-c_Lp$ would vanish and be positive immediately to its right, so its derivative there would be nonnegative, whereas \eqref{eq:numa-left-trap} shows that this derivative is strictly negative. Thus $c_Lp<h<sp$ on the part of the interval of existence to the left of $r_L$. Likewise, \eqref{eq:numa-right-trap} prevents a first downward crossing of $-tp$ to the right of $r_R$, giving $-tp<h<sp$ there.

These bounds also show that the maximal interval of existence is all of $(-1,1)$. Indeed, near a possible endpoint inside $(-1,1)$, the factors $x+h$ and $y-h$ stay bounded away from zero, so the graph of $h$ stays in a compact subset of the domain of $F_s$, and the continuation theorem for ordinary differential equations extends the solution beyond that endpoint. Hence the solution exists on $(-1,1)$ and satisfies the claimed bounds.
\end{proof}

\Needspace{7\baselineskip}
\begin{lemma}\label{lem:numa-thin-flat-extension}
Every solution furnished by \cref{lem:numa-thin-trapping}, extended by zero outside $(-1,1)$, is smooth on $\mathbb R$ and vanishes to infinite order at $\pm1$.
\end{lemma}
\begin{proof}
Let $h$ be one of the solutions in \cref{lem:numa-thin-trapping}. On either endpoint neighborhood put $\xi=h/p$; since $p>0$ on $(-1,1)$, this quotient is smooth there, and the trapping bounds place $\xi$ in a fixed compact subinterval of $(-1,1)$. Define the rational functions
\[
 R_L(p,\xi)=
 \frac{(1-s^2)(1+(1+s)p)}{(1+(1+\xi)p)(1-\xi)}-1,
\]
\[
 R_R(p,\xi)=1-
 \frac{(1-s^2)(1+(1-s)p)}{(1+\xi)(1+(1-\xi)p)}.
\]
The explicit formulas for the coefficients near the endpoints give
\[
 h'=q\,R_L(p,\xi)\quad\hbox{on the left},\qquad
 h'=q\,R_R(p,\xi)\quad\hbox{on the right}.
\]
Both rational functions are smooth, with all derivatives bounded, on a neighborhood of the relevant set of pairs $(p,\xi)$ with $p\geq0$ small and $\xi$ in the compact interval above. Writing $R_{L/R}$ for the function corresponding to the endpoint under consideration, we obtain from $h'=p'\xi+p\xi'$ and $q=p'$
\begin{equation}\label{eq:numa-scaled-ode}
 \xi'=\frac{p'}p\bigl(R_{L/R}(p,\xi)-\xi\bigr).
\end{equation}
For the chosen bump,
\[
 \frac{p'}p=-\frac{2r}{(1-r^2)^2}.
\]
This rational function and each of its derivatives grow at most like a power of the inverse distance to the endpoint. Moreover, each derivative of $p$ is $p$ times a rational function with the same property.

An induction using \eqref{eq:numa-scaled-ode} shows that every derivative $\xi^{(j)}$ grows at most polynomially in the inverse distance to the endpoint. The case $j=0$ is the trapping bound on $\xi$. For the induction step, differentiate the right-hand side of \eqref{eq:numa-scaled-ode}, including its term $-\xi$, $j$ times. Every resulting term is a product of a derivative of $p'/p$, a bounded partial derivative of $R_{L/R}$, and derivatives of $p$ and of $\xi$ of order at most $j$. All of these factors have polynomial bounds, by the induction hypothesis for the derivatives of $\xi$ and since every derivative of $p$ is bounded. This gives the bound for $\xi^{(j+1)}$.

Leibniz' rule for $h=p\xi$ now gives, for each $j$, constants $C_j,N_j$ such that
\[
 |h^{(j)}(r)|\leq C_j p(r)(1-|r|)^{-N_j}
\]
near the endpoints. The right-hand side tends to zero at the endpoints, because $p$ decays exponentially in the inverse distance to the endpoint. Thus all derivatives of $h$ tend to zero at $\pm1$. Extend each derivative by zero outside $(-1,1)$. By the mean value theorem, the difference quotient of $h^{(j)}$ at an endpoint tends to the limit of $h^{(j+1)}$, which is zero. Hence these extensions are the successive derivatives of the extension of $h$, which is therefore smooth and vanishes to infinite order at $\pm1$.
\end{proof}

\begin{proof}[of \cref{thm:numa-thin-zero-set}]
Let $\eta_0$ be as in \cref{lem:numa-thin-trapping}, and for $0<\eta<\eta_0$ denote by $h_\eta$ the solution with initial value $sp(0)-\eta$. Extend $h_\eta$ by zero, which gives a smooth function by the preceding lemma, and take a primitive $G_\eta$. It is constant near each end of the real line, so $G_\eta(r)$ extends smoothly to all of $\mathbb P^3$, including the two omitted projective lines. Set
\[
 \widetilde u_\eta=G_\eta(r)-\rho.
\]
On the open band,
\[
 \omega+\ddc\widetilde u_\eta
 =(x+h_\eta)P+(y-h_\eta)Q+(z+h_\eta')R.
\]
The coefficients of $P$ and $Q$ are strictly positive by \cref{lem:numa-thin-trapping}, while the equation gives
\[
 z+h_\eta'=\frac{k_s}{(x+h_\eta)(y-h_\eta)}>0.
\]
Outside the band the form is $\omega_0$, and on the two boundary hypersurfaces it is also the rank-one form $\omega_0$, since $h_\eta$ vanishes to infinite order there. It is therefore semipositive on $X$, and by \eqref{eq:numa-cube} its cube equals $6k_sP\wedge Q\wedge R=(\omega_0+sH)^3$ on the band and $\omega_0^3=(\omega_0+sH)^3$ elsewhere. Thus its cube is everywhere the measure $g\omega_X^3$ of \eqref{eq:numa-thin-density}. Finally, normalize by subtracting $\sup_X\widetilde u_\eta$.

Distinct values of $\eta$ give different slopes at $r=0$, so the corresponding potentials do not differ by a constant, and their normalized versions are distinct. This gives an uncountable family. The background is that of \cref{thm:degenerate-hermitian-nonuniqueness}, so it has bounded mass.
\end{proof}

\Needspace{8\baselineskip}
The density vanishes to infinite order on the two hypersurfaces. Their complement has three components, and on each of the two outer components, two different solutions differ by a constant, with different constants on the two components. The differential equation on the middle component joins these constants smoothly across the degenerate hypersurfaces. This is why ellipticity at every point of the positive-density set does not give a maximum principle argument across its boundary.

\section{Quadratic density zeros}
\label{sec:quadratic-density-nonuniqueness}

Infinite-order vanishing of the density is not needed for smooth nonuniqueness. We change the coefficients of $\omega_0$ so that its coefficient of $R$ stays positive and one of its coefficients of $P$ and $Q$ vanishes quadratically at each of the two hypersurfaces $M_\pm$. The resulting singular differential equation forces the slopes of differences of solutions to decay like $\exp(-A/d)$, where $d$ is the distance to the hypersurface and $A>0$. The slopes therefore extend smoothly by zero, even though the density vanishes only to finite order.

\Needspace{14\baselineskip}
\begin{theorem}\label{thm:quadratic-density-nonuniqueness}\index{Zero set!quadratic vanishing}
On $\mathbb P^3$ there exist a smooth nonclosed Hermitian form $\omega$ and a smooth nonnegative density $g$ such that:
\begin{enumerate}
\item The zero set of $g$ consists of two disjoint smooth real hypersurfaces. Near either hypersurface, $g$ is the square of a smooth defining function times a smooth strictly positive function.
\item The equation $(\omega+\ddc u)^3=g\omega_X^3$ has uncountably many distinct smooth $\omega$-plurisubharmonic solutions normalized by $\sup_Xu=0$.
\end{enumerate}
The background $\omega$ has bounded mass.
\end{theorem}

We keep the rank-one forms $P,Q,R$ and the function $r$ of \cref{subsec:numa-rank-directions}, and use new scalar coefficients.

Choose a smooth function $a$ valued in $[0,1]$, equal to one for $r\leq-2/3$ and zero for $r\geq-1/3$, and put $b(r)=a(-r)$. Choose a smooth even cutoff $\chi$ valued in $[0,1]$, zero for $|r|\leq5/4$ and one for $|r|\geq3/2$. Set $p(r)=(1-r^2)^2$ and
\begin{equation}\label{eq:quadratic-background}
 \begin{split}
 x&=(1-\chi)(p+a)+\chi\sigma,\\
 y&=(1-\chi)(p+b)+\chi(1-\sigma),\\
 z&=1-\chi+\chi\sigma',\qquad
 \omega_0=xP+yQ+zR.
 \end{split}
\end{equation}
Since these coefficients coincide with those of \eqref{eq:numa-fs} for $|r|\geq3/2$, the form $\omega_0$ extends by $\omega_X$ across the two omitted projective lines. The coefficient $z$ is strictly positive, and $x,y$ are strictly positive except that $y(-1)=0$ and $x(1)=0$. Thus $\omega_0$ is positive off $M_\pm=\{r=\pm1\}$ and has rank two there.

Choose a smooth cutoff $\kappa\in C_c^\infty((-2,2))$ equal to one on $[-3/2,3/2]$. For a parameter $\varepsilon>0$ to be chosen, define the correction
\[
 \rho=\varepsilon(r^2-1)\kappa(r),\qquad
 \omega=\omega_0+\ddc\rho.
\]
For sufficiently small $\varepsilon>0$, the form $\omega$ is strictly positive. Indeed, on $[-3/2,3/2]$ its coefficients are $x+2\varepsilon r$, $y-2\varepsilon r$ and $z+2\varepsilon$. The first is positive for $r\geq0$, and on $[-3/2,0]$ the function $x$ has a positive minimum, so it remains positive for small $\varepsilon$. Reflection gives the positivity of the second coefficient, and that of the third is immediate. On $3/2\leq|r|\leq2$, the form $\omega_0$ is strictly positive on a compact set and the correction is $O(\varepsilon)$. Near the projective ends there is no correction. The form $\omega$ is not closed: at $r=0$, $x'=y'=0$ and $z=1$, so by \eqref{eq:numa-differential}
\[
 \mathrm{d}\omega=\mathrm{d}r\wedge\bigl((x'-z)P+(y'+z)Q\bigr)\ne0.
\]

Define the density $g$ by
\[
 g\omega_X^3=\omega_0^3.
\]
It is smooth and positive exactly off $M_-\cup M_+$. By \eqref{eq:numa-fs} and \eqref{eq:numa-cube}, $g=xyz/(\sigma(r)(1-\sigma(r))\sigma'(r))$. Near either hypersurface, $xyz=(1+p)p$, and hence
\[
 g(r)=\frac{(1+p(r))(1-r^2)^2}{\sigma(r)(1-\sigma(r))\sigma'(r)}.
\]
The denominator is smooth and strictly positive at $\pm1$, and $1-r^2$ is a defining function for $M_\pm$ near these hypersurfaces. This proves the asserted quadratic vanishing. The potential $-\rho$ is one smooth solution.

\Needspace{15\baselineskip}
\begin{lemma}\label{lem:quadratic-density-trapping}
For every sufficiently small $\eta>0$, the equation
\begin{equation}\label{eq:quadratic-density-ode}
 h'=\frac{xyz}{(x+h)(y-h)}-z,\qquad h(0)=-\eta
\end{equation}
has a smooth solution on $(-1,1)$ with $-x<h<y$ and $h<0$. There are a constant $0<c<1$ and fixed neighborhoods of both endpoints, independent of $\eta$, on which
\[
 -cp<h<0.
\]
Every such solution, extended by zero, is a smooth function on $\mathbb R$ that vanishes to infinite order at $\pm1$.
\end{lemma}
\begin{proof}
The right-hand side $F(r,h)$ is smooth on $-x<h<y$, and $h=0$ is a solution. Near $-1$, $(x,y,z)=(1+p,p,1)$; near $1$, $(x,y,z)=(p,1+p,1)$. Fix $0<c<1$. For the trial function $-cp$, near the left endpoint
\[
 F(r,-cp)+cp'
 =\frac{1+p}{(1+(1-c)p)(1+c)}-1+cp'
 \longrightarrow-\frac{c}{1+c}<0,
\]
whereas near the right endpoint
\[
 F(r,-cp)+cp'
 =\frac{1+p}{(1-c)(1+(1+c)p)}-1+cp'
 \longrightarrow\frac{c}{1-c}>0.
\]
as $r\to\pm1$, since $p,p'\to0$ at both endpoints. The strict inequalities therefore hold on fixed neighborhoods of the endpoints, and $-cp$ is a lower barrier for the equation solved backwards near $-1$ and forwards near $1$.

Choose a compact interval $[r_L,r_R]\ni0$ with endpoints in those neighborhoods. Continuous dependence on initial values, around the solution $h=0$, gives $\eta_0>0$ such that for every $0<\eta<\eta_0$ the solution exists on $[r_L,r_R]$, stays in the domain of $F$, and lies above $-cp$ at $r_L$ and $r_R$. By uniqueness of solutions, it cannot meet the solution zero, so $h<0$. The barrier argument used in the proof of \cref{lem:numa-thin-trapping} then gives $-cp<h<0$ near both ends. On each compact subinterval of $(-1,1)$ the trapped graph stays at a positive distance from the boundary of the domain of $F$; hence, by the continuation theorem, the solution exists on all of $(-1,1)$.

\Needspace{6\baselineskip}
The equation factors as
\begin{equation}\label{eq:quadratic-density-logarithm}
 h'=h\,L(r,h),\qquad
 L(r,h)=\frac{z(x-y+h)}{(x+h)(y-h)}.
\end{equation}
Since $h<0$, we have $(\log(-h))'=L(r,h)$. Near the left endpoint, the trapping bounds and the explicit coefficients imply
\[
 L(r,h)=\frac{1+h}{(1+p+h)(p-h)}\geq\frac{C_1}{p}
\]
after shrinking that neighborhood, with $C_1>0$. Near the right endpoint they give
\[
 L(r,h)=\frac{-1+h}{(p+h)(1+p-h)}\leq-\frac{C_2}{p},
 \qquad C_2>0.
\]
Indeed, on the left, $1+h\geq1-cp\geq1/2$ and the denominator is at most $(1+p)(1+c)p$; on the right, the numerator is at most $-1$, and the denominator is at most $p(1+(1+c)p)$. Since $p$ is comparable to the square of the distance $d=1-|r|$ to the endpoint, integrating these inequalities backwards on the left and forwards on the right gives, for some $A>0$,
\begin{equation}\label{eq:quadratic-density-exponential}
 |h(r)|\leq C \mathrm{e}^{-A/d}
\end{equation}
near either endpoint. The constants may depend on the solution; uniformity in $\eta$ is not needed for the smooth extension.

It remains to control all derivatives. On the trapped graph, all partial derivatives in $(r,h)$ of the rational function $L$ grow at most polynomially in $d^{-1}$. Indeed, near an endpoint $x,y$ are polynomials in $r$, $z=1$, $h$ is bounded, and of the two factors in the denominator, one is bounded below by a positive constant and the other by $(1-c)p$. Repeated differentiation of the quotient therefore only introduces finite powers of $p^{-1}$.

Inductively, for every $j\geq0$ there are constants $C_j,N_j$ with
\[
 |h^{(j)}(r)|\leq C_j \mathrm{e}^{-A/d}d^{-N_j}.
\]
The case $j=0$ is \eqref{eq:quadratic-density-exponential}. Differentiate $h'=hL(r,h(r))$ $j$ times. Every resulting term is the product of a derivative of the factor $h$, of order at most $j$, and a derivative of the composite function $L(r,h(r))$. The latter involves only partial derivatives of $L$ and derivatives of $h$ of order at most $j$, which are polynomially bounded by the induction hypothesis, while the first factor has the exponential bound. This proves the case $j+1$. All derivatives of $h$ therefore tend to zero at $\pm1$. Extending them by zero and applying the mean value theorem to the difference quotients at the endpoints, as in \cref{lem:numa-thin-flat-extension}, shows that the extension of $h$ by zero is smooth and vanishes to infinite order at $\pm1$.
\end{proof}

\begin{proof}[of \cref{thm:quadratic-density-nonuniqueness}]
Let $\eta_0>0$ be such that \cref{lem:quadratic-density-trapping} applies for $0<\eta<\eta_0$, and let $h_\eta$ denote the solution with $h_\eta(0)=-\eta$. Take a primitive $G_\eta$ of its extension by zero. Since $G_\eta$ is constant near each end of the real line, $G_\eta(r)$ defines a smooth function on $X$. The potential $u_\eta=G_\eta(r)-\rho$ satisfies
\[
 \omega+\ddc u_\eta
 =(x+h_\eta)P+(y-h_\eta)Q+(z+h_\eta')R.
\]
Inside the band, the first two coefficients are positive by the lemma, and the equation gives
\[
 z+h_\eta'=\frac{xyz}{(x+h_\eta)(y-h_\eta)}>0.
\]
Outside the band the form is $\omega_0$, and on the two boundary hypersurfaces it is also the rank-two form $\omega_0$, since $h_\eta$ vanishes to infinite order there. Thus it is semipositive on $X$, and by \eqref{eq:numa-cube} its cube is $6xyzP\wedge Q\wedge R=\omega_0^3=g\omega_X^3$. Different initial values give potentials that do not differ by constants, so after normalization by the supremum we obtain an uncountable family. Bounded mass follows from \cref{prop:bmp-metric}, as for \cref{thm:degenerate-hermitian-nonuniqueness}.
\end{proof}

At each of the hypersurfaces $M_\pm$, the difference of two of these solutions agrees to infinite order with a constant, because its derivative is a difference of slopes that vanish to infinite order there. The constants at the two hypersurfaces need not agree. Thus the singular differential equation produces differences that are flat at $M_\pm$, even though the density has only quadratic zeros.

\section{Projective counterexamples with analytic density}
\label{uni-sec:examples}

The last construction prescribes the density before choosing the background. A symmetric pair of radial perturbations will have the same determinant, equal to a fixed polynomial in the global projective coordinate $t$. Patching the scalar coefficients while preserving their product makes the density real analytic on all of projective space. Its two quadratic zero hypersurfaces retain the separation that allows different extremal values.

\begin{theorem}\label{uni-thm:analytic-example}
Put
\[
 t=\frac{|A|^2}{|A|^2+|B|^2},\qquad
 a=\frac1{1+\mathrm{e}},\qquad b=\frac{\mathrm{e}}{1+\mathrm{e}},\qquad
 f_*(t)=(t-a)^2(t-b)^2.
\]
There exist a
smooth positive nonclosed $\omega$ and distinct smooth normalized
$\omega$-plurisubharmonic functions $u,v$ with
\begin{equation}\label{uni-eq:analytic-example}
 (\omega+\ddc u)^3=(\omega+\ddc v)^3=f_*(t)\omega_X^3.
\end{equation}
Both solution forms are positive away from the two connected real analytic
hypersurfaces $\{t=a\}$, $\{t=b\}$, and have complex rank two on them.
The density is globally real analytic and has precisely quadratic
transverse zeros.  For every $n\geq4$ the same assertions hold on
$\mathbb{P}^3\times\mathbb{P}^{n-3}$, with the pulled-back density and the product
Fubini--Study volume, and with rank $n-1$ on the two product
hypersurfaces.
\end{theorem}

\begin{proof}
The rank-one formula \eqref{eq:numa-cube} reduces the construction to scalar coefficients with a prescribed product. By \eqref{eq:numa-fs},
\begin{equation}\label{eq:analytic-fs-cube}
 \omega_X^3=6(\sigma')^2P\wedge Q\wedge R,
 \qquad \sigma'=\sigma(r)\sigma(-r)>0.
\end{equation}

\subsection{Prescribing the determinant}
\label{subsec:analytic-prescribed-determinant}

On $U$, one has $t=\sigma(r)$. We first construct two positive triples with product $f_*(\sigma)(\sigma')^2$ on the middle band. Their horizontal factors will lose exactly one direction at each endpoint.

Put $d_*(r)=f_*(\sigma(r))(\sigma'(r))^2$.  Since
$a=\sigma(-1)$, $b=\sigma(1)$, define the positive smooth divided
differences
\[
 k_-(r)=\int_0^1\sigma'(-1+s(r+1))\,\mathrm{d}s,\qquad
 k_+(r)=\int_0^1\sigma'(1+s(r-1))\,\mathrm{d}s,
\]
and $H_*(r)=k_-(r)k_+(r)\sigma'(r)>0$.  The fundamental theorem of
calculus gives
\begin{equation}\label{uni-eq:prescribed-det}
 \sigma(r)-a=(r+1)k_-(r),\quad
 \sigma(r)-b=(r-1)k_+(r),\quad
 d_*=(1-r^2)^2H_*^2.
\end{equation}
These formulas supply the smooth positive extension at both endpoints,
without a division by a vanishing function.  The symmetry
$\sigma(-r)=1-\sigma(r)$ makes $d_*$, and hence $H_*$, even.

Choose the flat bump
\[
 p(r)=
 \begin{cases}\mathrm{e}^{-1/(1-r^2)},&|r|<1,\\0,&|r|\geq1,
 \end{cases}
 \qquad q=p'.
\]
All derivatives vanish at $r=\pm1$: differentiating gives a finite sum
of powers of $(1-r^2)^{-1}$ multiplied by the exponential, and the
exponential dominates every such power.  On $|r|<1$,
$q=-2rp/(1-r^2)^2$.  Define
\begin{align}
 z_b&=\sqrt{q^2+2H_*^2},\qquad
 S_* =\sqrt{r^2+4p^2+2z_b(1-r^2)^2},\label{uni-eq:baseline1}\\
 x_b&=(S_*-r)/2,\qquad y_b=(S_*+r)/2.\label{uni-eq:baseline2}
\end{align}
Both radicands are strictly positive everywhere, so these are single
smooth formulas, including at the endpoints.  They satisfy
\begin{equation}\label{uni-eq:baseline-identities}
 z_b>|q|,\qquad x_b-y_b=-r,\qquad
 x_by_b=p^2+\frac{z_b}2(1-r^2)^2.
\end{equation}
The coefficients $x_b,y_b$ are nonnegative, positive away from the
endpoints, with values $(1,0)$ at $-1$ and $(0,1)$ at $1$.
On the open band,
\[
 S_*^2-(2p+|r|)^2=2(1-r^2)^2(z_b-|q|)>0,
\]
so $x_b>p$, $y_b>p$.  Consequently both triples
\begin{equation}\label{uni-eq:twotriples}
 (x_b-p,y_b+p,z_b-q),\qquad (x_b+p,y_b-p,z_b+q)
\end{equation}
are strictly positive there.

Here is the determinant cancellation.  Set
$D_*=x_by_b-p^2$, $E_*=p(x_b-y_b)=-rp$.
Equations \eqref{uni-eq:baseline-identities} and the formula for $q$ give
\[
 D_*=(z_b/2)(1-r^2)^2,\qquad qD_*=z_bE_*.
\]
The two products in \eqref{uni-eq:twotriples} therefore equal
\begin{equation}\label{uni-eq:det-cancellation}
 (D_*+E_*)(z_b-q)=(D_*-E_*)(z_b+q)
       =\frac{D_*}{z_b}(z_b^2-q^2)=d_*.
\end{equation}
For $|r|\geq1$, $p=q=0$ and the same identity reads
$x_by_bz_b=d_*$.  Both perturbed forms have exactly the prescribed
determinant, and rank two at the endpoints.

\subsection{Extension over the projective ends}
\label{subsec:analytic-global-patch}

The coefficients just constructed have the right determinant but need not extend over the projective ends. We interpolate the two horizontal coefficients to those of $f_*(t)^{1/3}\omega_X$ away from the zero levels, and determine the radial coefficient by the prescribed determinant.

Choose a smooth even cutoff $0\leq\eta\leq1$, equal to one on
$|r|\leq3/2$ and zero on $|r|\geq2$.  On $|r|<3/2$ set
$(x,y,z)=(x_b,y_b,z_b)$.  On $|r|>5/4$, where
$f_*(\sigma(r))>0$, put
\begin{align*}
 \kappa&=f_*(\sigma(r))^{1/3},\\
 x&=\eta x_b+(1-\eta)\kappa\sigma(r),\\
 y&=\eta y_b+(1-\eta)\kappa\sigma(-r),\qquad z=d_*/(xy).
\end{align*}
These coefficients are smooth and positive on the outer region.  On the
overlap $5/4<|r|<3/2$, $\eta=1$, $p=q=0$, and
$x_by_bz_b=d_*$; hence the two definitions agree identically.  For
$|r|\geq2$ they equal
$\kappa(\sigma(r),\sigma(-r),\sigma'(r))$.  Thus
\[
 \omega_0=xP+yQ+zR
\]
extends smoothly across both omitted projective lines, being equal there
to $f_*(t)^{1/3}\omega_X$.  This is smooth positive near those lines since
$f_*(0)=f_*(1)=a^2b^2>0$.  Globally $\omega_0$ is semipositive,
positive away from $r=\pm1$, and rank two on these levels.

The function
\[
 F(r)=2\int_0^r p(s)\,\mathrm{d}s
\]
is constant on each tail and extends to a smooth nonconstant function on
$X$.  Its half-Hessian is
$H_0=\ddc F/2=p(P-Q)+qR$.  By
\eqref{uni-eq:det-cancellation}, the forms
$A_0=\omega_0-H_0$, $B_0=\omega_0+H_0$ satisfy
\[
 A_0^3=B_0^3=f_*(t)\omega_X^3
\]
on the open block set and hence globally.  Their ranks and positivity
are those already checked in \eqref{uni-eq:twotriples}.  The construction is
symmetric: $x(-r)=y(r)$ and $z$ is even.

\subsection{A positive equation background}
\label{subsec:analytic-repair}

As in the first example, an exact correction makes the average form strictly positive without changing the two solution forms. Here the radial coefficient is already positive; the correction restores the missing horizontal direction at each endpoint. Choose a parameter $\varepsilon>0$ as restricted below.

Let $h(r)=(r^2-1)\eta(r)$, extend it by zero on the tails, and set
$\rho=\varepsilon h$.  The coefficients of
$\omega=\omega_0+\ddc\rho$ are
\begin{equation}\label{uni-eq:repair}
 (x+\varepsilon h',\ y-\varepsilon h',\ z+\varepsilon h'').
\end{equation}
On $|r|\leq3/2$, $h'=2r$, $h''=2$.
For $0\leq r\leq3/2$, $x+2\varepsilon r>0$: the only zero of
$x$ is at $r=1$, where the correction is positive.
On $[-3/2,0]$, $x$ has a positive minimum $m_0$, so positivity
holds if $3\varepsilon<m_0$.  Reflection proves positivity of the
second coefficient, and the third is $z+2\varepsilon>0$.
On the compact transition region $3/2\leq|r|\leq2$, all original
coefficients have a common positive lower bound $m_1$.  If
$M_1$ bounds $|h'|+|h''|$, taking also
$\varepsilon<m_1/(M_1+1)$ preserves their positivity there.
For $|r|\geq2$, the background equals
$f_*(t)^{1/3}\omega_X>0$.  This proves global strict positivity.

The correction is exact, so $\mathrm{d}\omega=\mathrm{d}\omega_0$.  At $r=0$,
evenness of $S_*$ gives $x'(0)=-1/2$, $y'(0)=1/2$.  Hence
\eqref{eq:numa-differential} gives
\[
 \mathrm{d}\omega_0=(z_b(0)+1/2)\,\mathrm{d}r\wedge(Q-P)\ne0.
\]
Indeed, if $\mathrm{d}r\wedge(Q-P)=0$, wedging with $P$ and $\dc r$
would contradict $P\wedge Q\wedge R>0$.  Thus the positive background
is nonclosed.

Now put $\widetilde u=-F/2-\rho$,
$\widetilde v=F/2-\rho$.  Then
$\omega+\ddc\widetilde u=A_0$,
$\omega+\ddc\widetilde v=B_0$.
Subtract their respective suprema to obtain $u,v$ with common
normalization zero.  Their difference remains nonconstant because $F$
is nonconstant.

\subsection{Zeros and higher dimensions}

It remains to check the global regularity of the prescribed density and the geometry of its zero levels. Taking a product with projective space then preserves the determinant equality and supplies examples in every higher dimension.

The function $t$ is real analytic in every projective chart, being a
quotient of real polynomials with positive denominator.  Hence $f_*(t)$
is globally real analytic.  Its two zero levels are regular because
$t=\sigma(r)$, $\sigma'>0$, and $\mathrm{d}r\ne0$.
Each level is connected: it is the quotient of the product of two
three-spheres of fixed nonzero radii by simultaneous circle multiplication.
It is compact, and a hypersurface has zero smooth volume.  The divided
differences in \eqref{uni-eq:prescribed-det} show near the two levels that
\begin{equation}\label{uni-eq:quadratic-zero}
 f_*(\sigma(r))=(r+1)^2g_-(r),\qquad
 f_*(\sigma(r))=(r-1)^2g_+(r),
\end{equation}
respectively, with smooth strictly positive transverse factors.

For $n\geq4$, let $\theta$ be the Fubini--Study form on
$Y=\mathbb{P}^{n-3}$, let $\pi_1:X\times Y\to X$ and $\pi_2:X\times Y\to Y$ be the projections, and set on $X\times Y$
\[
 \omega_n=\pi_1^*\omega+\pi_2^*\theta,\qquad
 \Omega_n=\pi_1^*\omega_X+\pi_2^*\theta,\qquad
 u_n=\pi_1^*u,\quad v_n=\pi_1^*v.
\]
Only the term with three factors from $X$ survives in an $n$-fold
product.  Consequently
\[
 (\omega_n+\ddc u_n)^n
 =\binom n3\pi_1^*(A_0^3)\wedge\pi_2^*(\theta^{n-3})
 =f_*(\pi_1^*t)\Omega_n^n,
\]
and the same holds for $v_n$.  The solution ranks on the two product
levels are $2+(n-3)=n-1$.  Restriction to a slice recovers
$\mathrm{d}\omega\ne0$, so $\omega_n$ is nonclosed.  Normalization,
distinctness, and real analyticity are preserved by pullback.
\end{proof}

Write $\operatorname{Vol}$ for volume with respect to $\omega_X^3$, or $\Omega_n^n$ in dimension $n$. The quadratic zeros complete the sharpness assertion of
Theorem~\ref{uni-thm:threshold}.  In finitely many tubular charts, the density
is bounded above and below by positive multiples of the square of a real
transverse coordinate, while the smooth volume has a positive bounded
Jacobian.  Away from these charts the density is bounded below.  Thus, with $s>0$ a sublevel parameter and $\alpha>0$ an integrability exponent,
\begin{equation}\label{uni-eq:sharp-asymptotic}
 \operatorname{Vol}\{f_*\leq s\}\asymp\sqrt s\quad(s\downarrow0),
 \qquad f_*^{-\alpha}\in L^1\quad\Longleftrightarrow\quad\alpha<\tfrac12
 \quad(\alpha>0).
\end{equation}
Thus every weaker inverse-power hypothesis with positive exponent below
$1/2$ admits smooth nonunique examples in every $n\geq3$.

\begin{remark}
For this example, $U=\{f_*(t)>0\}$ has three components: the middle band and the two outer regions. The solution forms have the same missing direction on each zero level, so $E_{\rm unif}=U$. The maximum and minimum of their potential difference occur on the two outer components, whose closures are disjoint, in accordance with \cref{uni-thm:anchors}. Although the middle component meets both outer closures, this adjacency does not propagate an extremal value through it.
\end{remark}


